\chapter{La representación y sus fisuras}


\label{cap:politica}
\section{El senador era el dueño de la finca}

Este capítulo describe superposiciones entre la propiedad del suelo y la representación
política del departamento. Conviene empezar por la primera en que el archivo muestra la banca ya ganada, porque
fija el patrón.

Desde mayo de 1927 el Boletín Oficial publica —seis viernes consecutivos, del 6 de mayo al 10 de junio— la nómina de
senadores electos de la provincia. Por el departamento de Caldera figura \textbf{Bernardo
Austerlitz\index{Austerlitz, Bernardo}}.

Dos años antes, en enero de 1925, tres pedimentos de cateo minero identificaban al
propietario del suelo sobre el que se solicitaba explorar. Dos de ellos —los expedientes
1164 C y 1165 C, sobre dos mil hectáreas cada uno en la finca San Alejo y Santa Rufina—
consignan que ese propietario es \textbf{Bernardo Austerlitz\index{Austerlitz, Bernardo}, domiciliado en la calle
Alberdi 374 de la ciudad de Salta}.

San Alejo es uno de los trece partidos policiales de 1909 y una de las fincas que el
artículo 6º del Plan de Desarrollo Urbano Ambiental sigue enumerando en 2015 como unidades
integrantes del municipio.

\begin{cautela}
Este libro \textbf{no afirma} que sean necesariamente la misma persona: el archivo entrega
un nombre completo poco frecuente en dos actos separados por dos años y referidos al mismo
departamento ---y en un tercero, anterior: en noviembre de 1914 un decreto lo nombra, junto a
Daniel Linares\index{Linares, Daniel}, jurado titular por Caldera para los reclamos de patentes (B.O.\ Nº 527)---, lo que es un indicio fuerte pero no una identificación registral. Y aun
siéndolo, ser propietario y ser senador por el mismo departamento no era, en 1927, ni
irregular ni infrecuente: el Senado provincial de la época se integraba en buena medida con
los apellidos de la propiedad rural, y la misma nómina lo muestra.

Lo que corresponde registrar es el patrón, porque el capítulo lo va a encontrar otra vez.
El propietario de una de las fincas mayores del departamento —una de las que la norma
todavía nombra— ocupa su banca en el Senado provincial. Casi un siglo después, este
libro documenta a un senador por el mismo departamento designado años antes gerente suplente
de una sociedad copropietaria del inmueble de un loteo.

Ninguno de los dos hechos prueba nada por sí solo, y este libro no los presenta como
equivalentes: median cien años, dos regímenes constitucionales y dos escalas económicas
distintas. Los pone juntos porque describen una continuidad que la tercera tesis enuncia
y que aquí se ve en su forma más simple: \textbf{en este departamento, quien tiene la
tierra ha tenido con regularidad también la representación}.
\end{cautela}

\section{Cuatro decretos en tres semanas}

Hay en 1928 una secuencia que conviene contar entera, porque reúne en tres semanas casi
todo lo que este libro discute: un derecho de agua, un conflicto entre municipios de aguas
arriba y aguas abajo, un funcionario que es parte interesada, y un cambio de gobierno que
revierte una decisión sin argumentar.

El punto de partida está en el capítulo \ref{cap:siglo}. El 28 de enero de 1924, la Provincia concedió a
\textbf{Francisco S.\ Urquiza\index{Urquiza, Francisco S.}} el uso de cien litros por segundo de los sobrantes del río
Wierna, para regar su finca Cerro de Buena Vista, a título precario y eventual. Entre los
fundamentos del decreto figuraba que, «oída la Comisión Municipal del expresado
departamento», ésta aconsejaba el despacho favorable.

\begin{ficha}{La secuencia de 1928}
\textbf{21 de abril.} Decreto 6877, gobernador Corbalán. \textit{Reclamación rechazada.}
Varios vecinos del \textbf{departamento de Campo Santo} y su Comisión Municipal habían
reclamado contra la concesión de Urquiza\index{Urquiza, Francisco S.}. El decreto no hace lugar, sobre la base de los
informes de la Dirección de Obras Públicas «evacuando los informes aconsejados por la
\textbf{Comisión Municipal de La Caldera} y el señor Fiscal General», que sostienen que la
acequia de Urquiza\index{Urquiza, Francisco S.} muere en el río Vaqueros y por lo tanto no perjudica a los reclamantes.
El derecho «seguirá ejerciéndose en la forma y condiciones acordadas».

\textbf{26 de marzo.} Decretos 6776 y 6777, uno detrás del otro: renuncian \textbf{el juez de paz
propietario Victorio Offredi y el comisario de policía Ramón Balvoa} ---así impreso---.
\textbf{Seis semanas antes de la acefalía y casi un mes antes del decreto que rechaza el
reclamo}, y ninguno de los dos da la causa.

\textbf{30 de abril.} Decreto 6958, exp.\ 674 U, mismo gobernador. Se acepta \textbf{la renuncia de
Francisco S.\ Urquiza\index{Urquiza, Francisco S.} al cargo de Miembro de la Comisión Municipal de La Caldera}. \textit{El decreto pasa del artículo 1º al 3º: no hay artículo 2 impreso.}

\textbf{9 de mayo.} Decreto 7096, ya con el gobernador Cornejo. Se declara en acefalía esa misma Comisión Municipal y se nombra comisionado \textit{ad honorem} a Luis E.\ Guardo; el 26 de junio el Decreto 9306 designa la nueva comisión: Julio Mercado, Manuel Condorí, Hermann Pfister\index{Pfister, Hermann}, Carlos Outes y Luis Leper.

\textbf{12 de mayo.} Decreto 8023, \textbf{gobernador Cornejo}. \textit{Sin efecto.}
«Visto el decreto \ldots\ por el que se concede al señor Francisco Urquiza\index{Urquiza, Francisco S.} el uso de los
sobrantes del agua del río Wierna \ldots\ y considerando: \textbf{que el mantenimiento de
dicha concesión es inconveniente}, el Gobernador de la Provincia decreta: Déjase sin
efecto la concesión \ldots»
\end{ficha}

\textit{El número y la fecha se cotejaron sobre el facsímil del Boletín Oficial Nº 1.222 (8 de junio de 1928, pág.~9), que publica el decreto con su refrendo ministerial. El salto de la numeración está en el original: la misma edición trae los decretos 7071 a 7099 del 8 y 9 de mayo y, desde ese mismo 9 de mayo, los 8000 en adelante. La edición siguiente (Nº 1.223) repite el patrón: del 8099, del 18 de mayo, pasa al 9000, del 19. La cifra de los miles avanza cada cien decretos, de modo que 7099--8000 y 8099--9000 son números consecutivos; con esa lectura el Decreto 9306 del 26 de junio encaja en la serie. Esa misma semana el nuevo gobierno renovó también al comisario de La Caldera (Decreto 7095, Julio Royo Ortiz, que el acto imprime «Ortíz»), al encargado de su Registro Civil (8002, Teófilo Reyes, en lugar de Ramón Balboa) y, el 18 de mayo, al receptor de rentas y expendedor de guías (8099, otra vez Royo Ortíz, en lugar de Balboa).}

\begin{cautela}
Hay que leer esto con cuidado, porque el archivo permite decir menos de lo que sugiere.

\textbf{Lo que está documentado.} Que en abril de 1928 el Estado rechazó un reclamo contra
la concesión de Urquiza\index{Urquiza, Francisco S.} apoyándose, entre otros, en el informe de la Comisión Municipal de
La Caldera. Que nueve días después Urquiza\index{Urquiza, Francisco S.} renunció a su cargo de miembro de esa comisión.
Que nueve días más tarde la comisión fue declarada en acefalía. Y que tres días después, con
otro gobernador, la concesión se revocó con una sola línea de fundamento.

\textbf{Lo que el Boletín de 1924
documenta.} Cabía preguntar si Urquiza\index{Urquiza, Francisco S.}
integraba la Comisión Municipal en 1924, cuando ese cuerpo aconsejó favorablemente su propia
solicitud. \textbf{El Boletín de ese año lo resuelve, y la respuesta es precisa}: el decreto
1.429 que concede el agua es del \textbf{28 de enero de 1924}, y el decreto \textbf{1.576 del 15
de abril} designa a Francisco Urquiza\index{Urquiza, Francisco S.} miembro de la Comisión
Municipal de La Caldera, junto con Julio Royo Ortiz, Camilo Cáceres y Serafín Yáñez (capítulo
\ref{cap:siglo}). \textbf{No integraba el cuerpo cuando el cuerpo informó: entró dos meses y
medio después.} La pregunta tiene respuesta, y la respuesta
favorece a Urquiza en ese punto; la secuencia, en cambio, queda más nítida que antes.

\textbf{Lo que sigue sin documentarse.} Si participó o se excusó en el informe de 1928, cuando
sí integraba la comisión. Y si la revocación de mayo
responde al conflicto con Campo Santo, al cambio de gobierno, a una revisión técnica o a
otra cosa: el decreto no lo dice, y una línea que declara «inconveniente» el mantenimiento
de un derecho no es un fundamento sino una fórmula.

La acefalía, en cambio, tiene un contexto que la vuelve menos elocuente de lo que parece: el mismo 9 de mayo, con la misma fórmula ---«\textit{para asegurar la existencia del régimen municipal}»--- y dos decretos antes ---el 7094---, el nuevo gobierno declaró en acefalía la Comisión Municipal de Embarcación. Fue parte de la reorganización que acompañó el cambio de gobernador, no una medida dirigida sólo contra La Caldera.

\textbf{Y el año entero agranda todavía más ese contexto.} Entre marzo y octubre de 1928 el departamento renovó \textbf{siete cargos}: juez de paz propietario, comisario, encargado del Registro Civil, receptor de rentas, subcomisario de Mojotoro, la comisión municipal entera y, en octubre, el juez de paz suplente. \textbf{No quedó uno solo en pie}, y el recambio empieza antes del episodio del agua (capítulo \ref{cap:siglo}). Eso no desmiente esta lectura: la encuadra.

Este libro no afirma, entonces, que haya habido conflicto de intereses ni represalia
política. Afirma que la secuencia ocurrió en el orden y en las fechas descriptas, y que
está publicada en el Boletín Oficial.
\end{cautela}

Hay tres cuestiones de esa secuencia que sobreviven hasta el presente del libro.

\textbf{El conflicto de aguas arriba y aguas abajo.} Quienes reclaman son los vecinos y la
comisión municipal de otro departamento, río abajo. Un siglo después, el capítulo \ref{cap:vaqueros}
documenta que el río Vaqueros separa dos jurisdicciones y que ninguna administra el
conjunto del cauce, y que la extracción de áridos se autoriza en ambas márgenes sin
autoridad común. El problema de la cuenca compartida está planteado en 1928 y todavía no
tiene solución institucional.

\textbf{La ausencia de mediciones como argumento disponible para cualquier lado.} El decreto
de 1924 otorgó precariamente porque no había aforos. El de abril de 1928 rechaza el reclamo
razonando sobre evaporación e infiltración, también sin aforos. Y el de mayo revoca sin
razonar en absoluto. Cuando nadie midió el río, cualquier conclusión es defendible.

\textbf{Y la superposición entre quien pide y quien informa.} Es el mismo rasgo que el capítulo \ref{cap:defensas} describe para 1910, cuando la comisión encargada de las defensas contra el
río estaba integrada por el presidente municipal, dos jurados y el comisario; y el mismo
que reaparece en el presente cada vez que el libro encuentra a la misma persona en dos
lados del expediente. En un departamento pequeño, la superposición no es una anomalía: es
la condición normal de funcionamiento, y por eso las reglas de excusación importan más ahí
que en cualquier otro lado.

\section{Y la Corte le dio la razón}

La secuencia de 1928 tuvo un final, y llegó tres años después.

Francisco S.\ Urquiza\index{Urquiza, Francisco S.} demandó a la Provincia. El 4 de marzo de 1931, la Corte de Justicia de
Salta, en ejercicio de la jurisdicción contencioso-administrativa que le confiere el
artículo 141 de la Constitución provincial, resolvió en su favor.

\begin{ficha}{Corte de Justicia de Salta, ``Urquiza\index{Urquiza, Francisco S.} c/ Gobierno de la Provincia'' --- 4 de marzo de 1931}
\textbf{Parte dispositiva:} «La Corte de Justicia \ldots\ \textbf{deja sin efecto el decreto
del P.\ E.\ de fecha Mayo 12 de 1928} impugnado por Francisco Urquiza\index{Urquiza, Francisco S.}, y, en consecuencia,
\textbf{ordena restablecer a éste en el ejercicio de la concesión de aguas} acordada según
decreto de fecha Enero 28 de 1924; sin costas \ldots\ y sin perjuicio a la acción ordinaria
que pueda corresponder al recurrente por los daños que alega.»

Firman los ministros Vicente Tamayo, David Saravia, Francisco F.\ Sosa, Humberto Cánepa y
Cristian Puló.

\textbf{Del voto del ministro Cánepa:} el Poder Ejecutivo «no debe revocar las concesiones
sino dentro de las cláusulas de las mismas o cuando se han vuelto contrarias a su finalidad
social». La concesión «no puede reputarse inconveniente por el solo hecho de subsistir». Y
tampoco es un mero permiso por haberse otorgado con carácter precario y eventual: «la
precariedad se refirió a la posible existencia de derechos anteriores y la eventualidad a la
posible existencia de sobrantes en la cantidad fijada».

\textbf{Del voto del ministro Saravia:} «los motivos de \textit{interés público} en que debe
basarse la revocación gubernativa \textbf{deben aparecer establecidos por formalidad y no
meramente invocados} por la resolución gubernativa; pues ello tendría la apariencia de una
arbitrariedad». Y agrega que la revocación de una concesión precaria sólo procede por
incumplimiento de condiciones o por superior interés público, «sujeto en este caso, claro
está, a una correspondiente indemnización».
\end{ficha}

Tres cosas hacen de este fallo una pieza importante para este libro, y ninguna tiene que ver
con quién ganó.

\textbf{La primera: fija qué significa «precario».} Durante casi un siglo, este libro
encuentra concesiones de agua otorgadas con ese carácter, y en el capítulo \ref{cap:loteo} documenta que
tres gestiones de abastecimiento poblacional recibieron tratamientos distintos sin que se
explicara por qué. La Corte, en 1931, dijo qué quería decir la palabra: la precariedad se
refiere a la posible existencia de derechos anteriores, y la eventualidad a la posible
existencia de sobrante. \textbf{No convierte al derecho en un permiso revocable a voluntad.}

\textbf{La segunda: fija un estándar de motivación.} Un acto que revoca un derecho no puede
limitarse a invocar el interés público: debe establecerlo. Un decreto cuyo único considerando
es «que el mantenimiento de dicha concesión es inconveniente» no cumple ese estándar, y la
Corte lo dice con todas las letras --- «tendría la apariencia de una arbitrariedad».

Ese estándar es exactamente el que este libro reclama, noventa años después, cuando el
capítulo \ref{cap:opacidad}, más adelante, muestra actos administrativos que publican la envoltura y reservan el criterio a
un anexo. \textbf{La exigencia de que el fundamento esté establecido y no sólo invocado no
es una demanda contemporánea ni una construcción de este trabajo: es doctrina de la Corte de
Justicia de Salta desde 1931.}

\textbf{Y la tercera: la ausencia de aforos vuelve a estar en el centro.} El fallo transcribe
el fundamento original de la precariedad, que es el mismo que el capítulo \ref{cap:siglo} documenta: no
existía aforo del río Wierna «ni dato positivo alguno sobre la condición actual de las aguas
con relación al uso que sobre ellas puedan tener terceros». La Corte no lo corrige ni lo
cuestiona: lo asume como el marco dentro del cual todos --- el concesionario, los vecinos de
aguas abajo, la Municipalidad y el propio Estado --- estuvieron discutiendo durante siete
años sobre un caudal que nadie había medido.

\begin{cautela}
Corresponde una precisión sobre el alcance de lo que este libro puede decir del episodio
completo. La Corte resolvió que la revocación fue infundada, y nada en su fallo se pronuncia
sobre lo que el capítulo describe unas páginas más arriba: que el informe municipal que
sustentó el rechazo del reclamo de Campo Santo provino de un cuerpo que el propio Urquiza\index{Urquiza, Francisco S.}
integraba.

Ese punto no fue materia del juicio. No lo alegó nadie, y la sentencia no lo menciona. De
modo que el fallo confirma que el Estado revocó mal, y no dice nada sobre si informó bien.
\end{cautela}

\section{Y lo tuvieron preso}

El fallo de 1931 no era el final. Un segundo pronunciamiento de la Corte, publicado en
septiembre de 1937, revela que entre medio hubo un proceso penal.

\begin{ficha}{Corte de Justicia, Primera Sala --- «Urquiza\index{Urquiza, Francisco S.} c/ Gobierno de la Provincia, daños y perjuicios»}
Sentencia del 28 de abril de 1936, que confirma con costas el rechazo de primera instancia
del 5 de octubre de 1935. Ministros: Humberto Cánepa, Vicente Tamayo y E.\ Cornejo Arias.

\textbf{Lo que el actor reclamaba:} \$20.000, de los cuales \$15.599,57 por daños materiales
y el resto por daño moral, derivados de la privación del uso del agua entre la revocación de
1928 y su anulación judicial.

\textbf{Lo que el fallo revela:} en agosto de 1928 --- cuatro meses después de renunciar a la
Comisión Municipal y de que se revocara su concesión --- Urquiza\index{Urquiza, Francisco S.} fue sometido a proceso penal
\textbf{«por hurto y uso indebido de agua»}, con \textbf{prisión preventiva}. La causa
terminó por \textbf{sobreseimiento definitivo por prescripción} en agosto de 1930.

\textbf{Quién lo denunció:} \textbf{Juan A.\ Bianchi\index{Bianchi, Juan A.}}, propietario de la finca «Vaqueros» o
«Entre Ríos».

\textbf{Por qué pierde:} la Corte razona que, si el proceso criminal fue promovido por
denuncia de un particular, no cabe responsabilizar al Estado por el daño moral; y que, aun
declarado ilegal el decreto de revocación, no corresponde indemnizar la privación del uso del
agua \textbf{sin que se demuestre que el actor habría podido usarla efectivamente} --- lo que
remite, otra vez, a que nadie sabía cuánta agua llevaba el río. Se descartan además la pericia
y las declaraciones testimoniales tomadas por el \textbf{Juez de Paz de La Caldera}, porque el
pliego de preguntas no obra en el expediente.

\textit{Las fechas vienen impresas con veinte años de más --- 1956 por 1936, 1955 por 1935,
1951 por 1931, 1950 por 1930 --- en un error sistemático del original o del escaneo. Se
corrigen por coherencia con el resto de la serie y se advierte aquí.}
\end{ficha}

\textbf{La primera es que el episodio de 1928 fue mucho más grave de lo que el archivo
administrativo dejaba ver.} Los decretos publicados mostraban un reclamo rechazado, una
renuncia, una acefalía y una revocación. Faltaba lo principal: \textbf{el concesionario fue
detenido}. Un conflicto por sobrantes de agua entre dos fincas vecinas terminó con uno de los
propietarios en prisión preventiva y con una causa penal que sólo se cerró por prescripción,
dos años después.

\textbf{La segunda es quién denunció.} No el Estado ni el municipio de aguas abajo: el dueño
de la finca lindera, \textbf{la misma «Vaqueros o Entre Ríos» que en 1935 sería rematada por
el Banco Provincial} y cuyo lindero este era, según ese edicto, el propio Urquiza\index{Urquiza, Francisco S.}; para entonces el ejecutado era Ramón Bracheri, de modo que la finca cambió de titular entre 1928 y 1935.

De modo que la secuencia completa es: dos fincas contiguas, un río, una concesión de
sobrantes otorgada a una de ellas porque nadie había medido el caudal, una revocación administrativa
sin fundamento, una denuncia penal del vecino, una prisión preventiva, un fallo que
anula la revocación, una demanda de daños y un rechazo definitivo. \textbf{Doce años entre la concesión de 1924
y el fallo de 1936 por agua que nadie sabía si existía.}

\textbf{Y la tercera es que el Estado no pagó, y por qué.} La Corte había declarado ilegal la
revocación en 1931, pero en 1936 rechaza la indemnización porque el actor no demostró que
hubiera podido usar esa agua. \textbf{La misma ausencia de aforos que hizo precaria la
concesión en 1924 le impidió cobrar los daños en 1936.} El desconocimiento del recurso operó,
sucesivamente, a favor del Estado en las dos instancias.

\begin{cautela}
Corresponde consignar una discrepancia, y quedó resuelta. Este fallo fecha la revocación el \textbf{12 de marzo
de 1928}; el decreto 8023, publicado en el Boletín Nº 1.222 del 8 de junio de 1928, lleva fecha
\textbf{12 de mayo}, y el fallo de 1931 anula expresamente «el decreto del P.\ E.\ de fecha Mayo 12 de 1928». Además, el 12 de marzo gobernaba todavía Corbalán, que el 21 de abril ratificó la concesión. El error es del fallo de 1936, el mismo que imprime las fechas con veinte años de más.

Y corresponde una precisión sobre lo que este libro no puede decir. Se ignora si el juicio
por hurto de agua tuvo o no fundamento: el sobreseimiento fue por prescripción, no por
inocencia, y el fallo no describe los hechos imputados. Que la causa terminara sin condena no
prueba que no hubiera habido uso indebido, ni lo contrario.
\end{cautela}

\section{Títulos que hay que reponer}

Hay un segundo rasgo del archivo temprano que conviene traer aquí, porque explica por qué
la propiedad de este territorio es tan difícil de reconstruir.

En septiembre de 1926, el Boletín publica una citación en los autos «sobre reposición de
título de propiedad de la finca denominada Peñones o Los Peñones, ubicada en el
Departamento de la Caldera», promovido por el doctor Francisco F. Sosa en representación
de la señorita María Isabel Gómez.

\textbf{Reponer un título es reconstruir judicialmente un instrumento que se perdió.} Los
Peñones es, además, otro de los partidos de 1909 —figuraba como «Peñones y Chalchanio»,
con Liborio Guerrero\index{Guerrero, Liborio} como comisario auxiliar.

De modo que en las dos primeras décadas del Boletín el departamento registra: un pleito de
1908 que se resuelve, por prescripción treintañal, porque el título de compra no figura en autos; una mensura de 1910 cuyos límites se describen por
«unos mojones antiguos»; y en 1926 un juicio para reponer el título perdido de otra finca.

La opacidad dominial de este territorio no es una consecuencia del mercado inmobiliario
contemporáneo ni de la digitalización incompleta de los registros. Es un rasgo antiguo y
documentado, y los expedientes que lo prueban están publicados.

\subsection*{Tres cargos, una aritmética}

El departamento elige un senador y un diputado provinciales, y cada municipio su intendente. Tres de esos cuatro cargos ---la banca del Senado y las dos intendencias--- producen, cada cuatro años, una configuración de tres actores que puede alinearse o enfrentarse, y cuya composición determina la capacidad de negociación del departamento frente al Ejecutivo provincial.

\section{Ciclo 2021}

La banca senatorial fue disputada por la senadora en ejercicio María Silvina Abilés y por el ex intendente y ex diputado Miguel Calabró. Los dos intendentes del departamento, Daniel Moreno por Vaqueros y Diego Sumbay\index{Sumbay, Diego Ramiro} por La Caldera, acompañaron, según la cobertura de prensa, la candidatura oficialista; Moreno, además presidente del Foro de Intendentes, encabezó la lista de convencionales constituyentes del mismo frente.

\begin{ficha}{María Silvina Abilés --- Cámara de Senadores de Salta}
\begin{itemize}[leftmargin=*,nosep]
\item Senadora \textbf{provincial} por el departamento La Caldera. Mandatos 2013--2017 y 2017--2021.
\item Vicepresidente Tercero de la Cámara durante el segundo mandato.
\item Partido Memoria y Movilización; bloque Salta Tiene Futuro.
\item No fue reelecta en agosto de 2021; cesó el 18 de noviembre de ese año, consignando en su discurso más de quinientos proyectos presentados.
\item Actividad departamental documentada en 2021: cámaras de videovigilancia para todo el departamento (Expte. 90-30.016/2021); refuncionalización del Colegio Secundario Senado Provincial de La Caldera (90-30.017/2021); puente peatonal en el ingreso de Villa Sara (90-30.027/2021); adoquinado de la calle Catalán, Vaqueros (90-30.028/2021); obras en la Escuela Nº 136 Juana Moro de López.
\end{itemize}
\end{ficha}

\begin{cautela}
La calificación de ``ex senadora nacional'' que consignó parte de la cobertura de enero de 2026 es incorrecta, y queda corregida conforme el registro oficial de la Cámara.

Y una nota lateral que conviene no perder: el expediente 90-30.016/2021, sobre cámaras de videovigilancia para todo el departamento, es un punto de entrada documentado a H.6 del libro anterior --- la hipótesis de que la inversión provincial en sistemas de vigilancia urbana crece con transparencia limitada y con potencial de uso superior al objetivo declarado --- ahora en escala departamental.
\end{cautela}

\subsection*{Ciclo 2025}

En 2025 la configuración se invirtió. La interna oficialista enfrentó al senador en ejercicio Héctor Miguel Calabró\index{Calabró, Héctor Miguel} contra el intendente de Vaqueros, Daniel Roberto Moreno Ovalle\index{Moreno Ovalle, Daniel Roberto}, con apoyo del intendente de La Caldera; Calabró contó con el respaldo del diputado Luis Gerardo Mendaña\index{Mendaña, Luis Gerardo}. La cobertura provincial consignó que el gobierno pidió unificar la lista y que el pedido no fue atendido. El escrutinio definitivo del Tribunal Electoral dio a Moreno \textbf{2.989 votos}, el 27,94 por ciento de los votos positivos, contra \textbf{2.830} de Calabró, el 26,45; La Libertad Avanza quedó tercera con 2.026. La prensa había informado 2.988 votos y una participación superior al 65 por ciento. Moreno dejó el Ejecutivo municipal para asumir la banca el 10 de diciembre de 2025. Compitieron además, según la prensa, candidatos de dos líneas libertarias --- entre ellos un ex oficialista que pasó a La Libertad Avanza en 2023 y fue designado al frente del organismo previsional nacional en la provincia --- y de fuerzas locales menores.

La elección se decidió en Vaqueros. En el municipio de La Caldera ganó Calabró, 1.531 votos contra 988; en Vaqueros ganó Moreno, 2.001 contra 1.299. El intendente del municipio más poblado le ganó la banca al senador que se impuso en el municipio que da nombre al departamento.

\begin{ficha}{Escrutinio definitivo, elecciones provinciales del 11 de mayo de 2025 --- categoría senador, departamento La Caldera}
\begin{itemize}[leftmargin=*,nosep]
\item \textbf{Mesas:} cincuenta y una, en nueve establecimientos: cuatro en el municipio de La Caldera y cinco en el de Vaqueros.
\item \textbf{Votos computados:} \textbf{11.171}, de los cuales 10.698 positivos y 473 en blanco. La planilla no consigna votos nulos.
\item \textbf{Electores habilitados:} \textbf{17.068}, según la planilla oficial de establecimientos y mesas.
\item \textbf{Por municipio:} \textbf{3.697 votos en La Caldera} y \textbf{7.474 en Vaqueros}.
\item \textbf{Promedio por mesa:} 219 votos.
\end{itemize}
\end{ficha}

El Tribunal Electoral publica también el padrón, mesa por mesa. \textbf{En mayo de 2025 el departamento tenía 17.068 electores habilitados}: 5.379 en el municipio de La Caldera y 11.689 en el de Vaqueros. Los 11.171 votos computados son el 65,4 por ciento de ese padrón, sin contar los nulos, lo que confirma la participación que informó la prensa.

\needspace{10\baselineskip}
{\sffamily\bfseries\footnotesize Electores habilitados en el departamento La Caldera --- Tribunal Electoral de Salta\par}
\vspace{0.5em}
{\small
\setlength{\LTpre}{0.6em}\setlength{\LTpost}{1.2em}
\begin{longtable}{@{}p{4.2cm}rrr@{}}
\toprule
\textbf{Elección} & \textbf{La Caldera} & \textbf{Vaqueros} & \textbf{Departamento} \\
\midrule
\endhead
Primarias, agosto de 2017 & 3.700 & 8.599 & 12.299 \\
Generales, noviembre de 2019 & 4.056 & 9.309 & 13.365 \\
Generales, agosto de 2021 & 4.339 & 9.738 & 14.077 \\
Generales, mayo de 2023 & 4.956 & 10.816 & 15.772 \\
\textbf{Generales, mayo de 2025} & \textbf{5.379} & \textbf{11.689} & \textbf{17.068} \\
\bottomrule
\end{longtable}
}

\begin{figure}[htbp]
\centering
\includegraphics[width=\textwidth]{fig-padron-poblacion.png}
\caption[Padrón electoral y población del departamento, 2017--2025]{\textbf{Padrón electoral y población del departamento La Caldera, 2017--2025.} Electores habilitados por municipio en cada elección provincial, frente a la población censada en 2010 y 2022 y a los mayores de dieciocho años de 2022. La barra de 2017 va rayada: su total coincide dígito por dígito con la población de 2022 y puede ser un error de carga. Elaboración propia con datos del Tribunal Electoral de Salta y del INDEC, con \texttt{scripts/grafico\_padron.py} sobre \texttt{datos/electoral/} (repositorio \texttt{mapas\_caldera}).}\label{fig:padron}
\end{figure}

\begin{cautela}
Las cifras oficiales confirman el problema que la prensa dejaba entrever, y lo precisan. La lámina \ref{fig:padron} lo muestra en una sola imagen.

\textbf{El padrón es casi el doble de la población adulta.} El Censo 2022 registró en el departamento 12.299 habitantes, de los cuales 8.976 mayores de dieciocho años. El padrón de 2025 tiene 17.068 electores: \textbf{el 139 por ciento de toda la población censada}, niños incluidos. En la provincia, la misma relación es del 76 por ciento ---1.092.561 electores sobre 1.441.351 habitantes---, que es lo esperable en una población joven. \textbf{El departamento está más de sesenta puntos por encima de la media provincial}, y en sentido contrario al esperable.

\textbf{La diferencia es mayor en Vaqueros.} Su padrón equivale al 150 por ciento de sus 7.774 habitantes; el de La Caldera, al 119 por ciento de sus 4.525.

\textbf{Y no es reciente.} Según la planilla, en 2017 el departamento tenía 12.299 electores \textit{---cifra idéntica, dígito por dígito, a la población censal de 2022, lo que sugiere un error de carga---}, muy por encima de los 7.763 habitantes del censo de 2010. \pendiente{cotejo de la planilla de electores de 2017 contra el original del Tribunal Electoral} Si la cifra se confirma, entre 2017 y 2025 el padrón creció un 38,8 por ciento, a un ritmo del 4,2 por ciento anual, algo mayor que el 3,9 anual con que creció la población entre 2010 y 2022.

Este libro \textbf{no afirma} irregularidad alguna. La explicación más probable es legal y conocida: se vota donde figura el domicilio del documento, no donde se vive, y el Censo 2022 ---el primero \textit{de derecho} del país--- cuenta a las personas donde residen habitualmente, que para el dueño de una casa de fin de semana no es el departamento. Un departamento con casi una de cada cinco viviendas sin residentes permanentes, clubes de campo y casas de fin de semana de vecinos de la capital (capítulo \ref{cap:poblacion}) es exactamente el lugar donde cabe esperar muchos electores que no son habitantes. También puede influir que el padrón tarde en dar de baja a quienes se mudaron.

Si es así, el dato es de primer orden para este libro: \textbf{entre dos quintos y casi la mitad de quienes pueden elegir a las autoridades del departamento no vive en él} ---según se tome la población adulta de 2022 o se la proyecte a 2025, y sin descontar a quienes el padrón no dio de baja---, y es probablemente la misma población ---propietarios no residentes--- que el capítulo \ref{cap:redes} identifica entre los beneficiarios de la falta de infraestructura. Lo que falta para confirmarlo no es el tamaño del padrón, que ya se conoce, sino su composición: cuántos electores se incorporaron por cambio de domicilio y desde dónde, dato que el Tribunal Electoral no publica.
\end{cautela}




\section{1936: una mesa anulada}

Diez años antes de la convocatoria de 1946, que la sección siguiente transcribe, hay un episodio que este relevamiento no buscaba y que
conviene dejar registrado, porque cruza las dos series que el libro sigue por separado: la
electoral y la de los nombres.

\begin{ficha}{Nómina de presidentes de comicio insaculados para 1936--1937 --- Colegio Electoral Nº 2, La Caldera}
\textbf{Circuito 5, Mesa 1} (Escuela Elemental Mixta): presidente \textbf{Regis,
Augusto}\index{Regis, Augusto}; suplente 1º Robles, Eusebio José; suplente 2º \textbf{Linares,
Daniel}\index{Linares, Daniel}.

\textbf{Circuito 5, Mesa 2} (Iglesia Parroquial): presidente Alonso, Fabián; suplentes Romano,
Lorenzo, y Lozano, Juan de Dios.

\textbf{Circuito 6, Mesa 1} (Registro Civil de Vaqueros): presidente Bruzzo, Atilio José.

El encabezado del documento aclara el método: los presidentes de mesa eran
\textbf{insaculados} ---sorteados--- por el Tribunal Electoral.
\end{ficha}

Es la única pieza del archivo que pone a Augusto Regis y a Daniel Linares en la misma mesa
electoral. No es la primera que los reúne: en 1910 integraron la misma comisión de defensas, y
en 1922 un mismo decreto los designó miembros de la Comisión Municipal. El capítulo \ref{cap:fincas} los sigue por separado durante treinta años
---presidente municipal y propietario el primero, presidente de la comisión de defensas de
1910 y propietario de tres fincas el segundo---; aquí están, en 1936, como presidente y
suplente segundo de la misma mesa. \textbf{Y por sorteo}, lo que en un padrón de pocos
centenares de electores dice menos sobre la voluntad de nadie que sobre el tamaño del conjunto
del que se sorteaba.

\textbf{Esa mesa fue anulada.} Por decreto del 6 de marzo de 1936, el Ministerio de Gobierno
transcribe la comunicación del Tribunal Electoral ---firmada por David Saravia y Pedro J.\
Aranda--- que informa haber anulado la elección del domingo 1º de marzo en una lista de mesas,
entre ellas la \textbf{Mesa Nº 1 del Circuito Nº 5 de La Caldera}, y convoca a comicios
complementarios para el \textbf{domingo 15 de marzo}, por el artículo 77 de la Ley 122. Se
elegían gobernador, vicegobernador y representantes a la Legislatura.

\begin{cautela}
El decreto \textbf{no dice por qué} se anuló ninguna de las mesas de la lista: transcribe el
resultado del examen de urnas y documentación, no su fundamento. Este libro no afirma
irregularidad alguna ni la vincula con quién presidía: la anulación de mesas por vicios
formales ---documentación incompleta, actas mal labradas--- era y es corriente, y La Caldera
figura en una lista junto a mesas de otros departamentos. Lo que queda establecido es el hecho
y su fecha.
\end{cautela}

\section{Y en 1946 ya elegía uno}

Antes de llegar a 1983 conviene poner una referencia mucho más antigua, porque muestra que el
tamaño de la representación departamental no es un efecto de la reforma de 1986 ni del
crecimiento reciente.

\begin{ficha}{Decreto Nº 9686-G --- Salta, 15 de diciembre de 1945 --- Expte.\ Nº 2949/945}
Convoca al pueblo de la Provincia a elecciones nacionales y provinciales simultáneas para el
\textbf{domingo 24 de febrero de 1946}, con sujeción a los arts.\ 27, 70 y 71 de la Ley 122 de
Elecciones de la Provincia. Dictado en Acuerdo de Ministros por el Cnel.\ Ángel W.\ Escalada.

Su art.\ 4º reparte la representación legislativa en \textbf{tres escalones}:

\begin{itemize}[leftmargin=*,itemsep=1pt]
\item \textbf{La Capital}: siete diputados titulares y siete suplentes, un senador titular y uno suplente.
\item \textbf{Rosario de la Frontera, Rosario de Lerma, Metán, Anta y Orán}: dos diputados titulares y dos suplentes, un senador titular y uno suplente cada uno.
\item \textbf{Los dieciséis restantes} ---Rivadavia, La Candelaria, Iruya, Santa Victoria, \textbf{La Caldera}, Campo Santo, Chicoana, Cerrillos, La Viña, Guachipas, Cafayate, Molinos, San Carlos, Cachi, La Poma y Los Andes---: \textbf{un} diputado titular y uno suplente, \textbf{un} senador titular y uno suplente.
\end{itemize}

Publicado en diez ediciones consecutivas del 19 al 31 de diciembre de 1945 y en otras cuarenta y
cinco, del 2 de enero al 23 de febrero de 1946 (Nº 2.465 a 2.509), víspera del comicio, como manda
su propio art.\ 13º: cincuenta y cinco ediciones en total, según su marca de publicación
«e|19|12|45 -- v|23|2|46».
\end{ficha}

El comicio se hizo. Por La Caldera resultaron electos el senador \textbf{José Luis Alvarez} y el diputado \textbf{Darío Arias}, a quienes el Decreto 11198-G del 22 de abril de 1946 convoca, con los demás legisladores, a constituir las Cámaras el 26 de abril (B.O.\ Nº 2.554). La Ley 745, sancionada el 29 de agosto, fijó en esa fecha el comienzo del mandato y repartió los departamentos en dos grupos para sortear quiénes se renovaban a los dos años y quiénes a los cuatro; La Caldera ---«Caldera», sin artículo--- quedó en el primero (B.O.\ Nº 2.652). El Boletín de 1946 no publica el resultado del sorteo, que la ley manda hacer en el primer período de sesiones ordinarias.

Ochenta años después, y con el departamento cuadruplicado ---2.831 habitantes en el censo de 1947, 12.299 en el de 2022---, \textbf{la asignación es la misma}:
un diputado y un senador. \textbf{Y en 1949 la convocatoria incluye el municipio}: el decreto 17.056-G del 22 de septiembre, que llama a elecciones para el 27 de noviembre, le asigna a La Caldera, además del senador y el diputado, \textbf{un intendente y tres concejales} titulares, con sus suplentes (B.O.\ Nº 3514, h.\ 23). Si se votó y quién resultó electo no está en el corpus de 1949; en abril de 1950 el decreto 1406-G nombra jueces de paz a propuesta del «\textit{señor Intendente de la Municipalidad de La Caldera}», sin decir quién es, y ningún acto de 1950 da su nombre (B.O.\ Nº 3683, h.\ 4). El de 1951 lo da de pasada: entre los proveedores de unas cañerías figura «\textit{Gregorio M.\ García, Intendente de La Caldera}» (Res.\ 953-A, Nº 3987, h.\ 7), y en septiembre la convocatoria para el 11 de noviembre vuelve a asignarle al municipio un intendente y tres concejales (8298-G, Nº 4032, h.\ 18 a 20), con un resultado que el Boletín del año tampoco publica; en 1952 el intendente preside la comisión del hogar de niños del pueblo y ningún acto del año lo nombra (13108-A, Nº 4201, h.\ 6), y la Ley 1433 sortea al departamento en el grupo A para la renovación de las Cámaras (Nº 4202, h.\ 4); en 1953 la Municipalidad propone la terna para el juez de paz suplente (4683-G, Nº 4413, h.\ 9) y ningún acto del año nombra al intendente; en 1954 el decreto 8545-G convoca para el 25 de abril a elegir un intendente y tres concejales para 1955--1958 (Nº 4609, h.\ 7), y ningún acto del año nombra al intendente, ni al que estaba ni al que resultó electo; en 1955 el interventor federal declara intervenida la Municipalidad para «\textit{una mejor organización dentro del régimen municipal}» y le nombra un interventor, sin nombrar al intendente saliente (247-G, Nº 5032, h.\ 5); en 1956 el interventor renuncia y se designa otro (4488-G y 4489-G, Nº 5253, h.\ 9), y en 1957 la convocatoria para el 23 de febrero de 1958 le asigna al municipio \textbf{cuatro miembros titulares y cuatro suplentes de la comisión municipal}, y ningún intendente (11.420-G, Nº 5543, h.\ 4 y 5). En 1983, como se ve más adelante, la boleta tampoco traía intendente. Lo que cambió en el intervalo ---la reforma de 1986, la elección
directa del intendente, el crecimiento demográfico más alto de la provincia--- no tocó el
escalón en el que La Caldera está colocada dentro del reparto legislativo.

\section{Qué era La Caldera en 1983}

La búsqueda de los primeros intendentes del retorno democrático encontró otra cosa, mejor: encontró que no había intendentes. Cuatro ediciones del Boletín Oficial de 1983 permiten reconstruir qué era, institucionalmente, el municipio en el momento en que empieza el período que este libro estudia. \textbf{Y en 1983 no era nuevo}: tras la convocatoria de 1957, que le dio al municipio comisión municipal y no intendente, un decreto de mayo de 1958 designa a Werfil Gallo entre «\textit{los señores Intendentes Municipales}» encargados de las comisiones viales (391-E, Nº 5667, h.\ 19), y en 1959 y 1960 el Poder Ejecutivo lo designa \textbf{presidente de la comisión municipal} por un año más, porque el artículo 182 de la Constitución dispone que esos presidentes «\textit{durarán un año en el desempeño de sus funciones, pudiendo los mismos ser designados nuevamente}» (6726-G, Nº 5908, h.\ 10; 15.649-G, Nº 6283, h.\ 12); en marzo de 1960 el municipio vuelve a elegir sólo los cuatro miembros titulares y cuatro suplentes de la comisión (10324-G, Nº 6051, h.\ 6). La comisión se votaba; a su presidente lo designaba la Provincia. \textbf{Y en 1961 y 1962 la designación cambia de nombre}: en junio de 1961 un decreto vuelve a designar a Gallo por un año, con el mismo artículo 182 (18308-G, Nº 6411, h.\ 8); en noviembre la intervención federal declara en comisión a las autoridades municipales de toda la provincia; en febrero de 1962 la Intervención designa a Gallo comisionado interventor comunal de La Caldera (Nº 6563, h.\ 8), y en mayo otro decreto le da por terminadas las funciones y pone en su lugar a Augusto Regis ---el nombre del que presidió la comisión municipal en los años treinta y cuarenta; el acto no dice si es el mismo--- (2382-G, Nº 6615, h.\ 11). Entre las dos designaciones, la Intervención había convocado para el 18 de marzo a elegir cuatro miembros titulares y cuatro suplentes de la comisión municipal (896-G, Nº 6547, h.\ 3 a 5), y ningún acto hallado de 1962 da el resultado. \textbf{En 1963}, también leído con pendientes, Regis renuncia desde el 28 de febrero y la Intervención designa comisionado interventor a Julio Catalán Arrellano ---«Arellano» en el acto de octubre; el nombre del primer presidente de la comisión municipal de Vaqueros, designado en 1970, sin que conste que sea el mismo--- (6751-G, Nº 6823, h.\ 9); la convocatoria a elecciones de enero vuelve a darle al municipio cuatro miembros titulares y cuatro suplentes de la comisión municipal (decreto-ley 236-G, Nº 6793, h.\ 7), y ningún acto hallado del año da el resultado; y el 15 de octubre, con el gobierno constitucional que asumió el 12, un decreto da por terminadas las funciones de los comisionados interventores de veintitrés localidades y designa presidente de la comisión municipal de La Caldera a Mario García (decreto 32, Nº 6965, h.\ 5 y 6). \textbf{En 1964 y 1965}, leídos con pendientes, aparece en los actos de la justicia de paz un cuerpo que la convocatoria no elige: en junio de 1964 la Provincia designa al juez de paz titular y al suplente por la terna que eleva el «\textit{H.\ Consejo Deliberante de la Municipalidad de La Caldera}» (3550, Nº 7118, h.\ 7 y 8), y en octubre de 1965 al titular, por la que aprueba el «\textit{H.\ Concejo Deliberante de la Municipalidad}» (10619, Nº 7453, h.\ 7), mientras la convocatoria para el 14 de marzo de 1965 le asigna al municipio, otra vez, cuatro miembros titulares y cuatro suplentes de la comisión municipal y ningún concejal (6454, Nº 7243, h.\ 12); qué es ese concejo los actos no lo dicen. En abril de 1965 renuncia el presidente de la comisión municipal, Gregorio M.\ García ---el nombre del intendente de 1951; el presidente designado en 1963 figura como Mario García, y que sean la misma persona no consta---, porque fue electo senador provincial en esos comicios (8331, Nº 7336, h.\ 6), y el 14 de mayo la Provincia designa presidente, invocando el artículo 178 de la Constitución, a Esteban Mogro, a quien una semana antes se le había aceptado la renuncia al juzgado de paz (8560, Nº 7347, h.\ 18; 8430, Nº 7343, h.\ 9). \textbf{En 1966}, también con pendientes, el Boletín cambia tres veces de gobierno ---el constitucional hasta fines de junio, una intervención federal en julio, con la que la numeración de los decretos vuelve a empezar, y desde agosto el general Héctor D'Andrea como gobernador---, y el 7 de septiembre el decreto 777 da por terminadas las funciones de Mogro como presidente de la comisión municipal, con las de las autoridades de otros cinco municipios, y designa en su lugar a Julio González (Nº 7666, h.\ 7); ese mes renuncia el juez de paz titular designado en 1965, y ningún acto hallado del año cubre la vacante (973, Nº 7672, h.\ 9). \textbf{En 1967}, también con pendientes, la vacante se cubre en febrero, por terna de la Municipalidad, con Andrés López Arrieta como titular y Mogro como suplente (3069, Nº 7774, h.\ 12 y 13), y el titular se jubila desde el 31 de marzo sin que lo hallado del año le nombre sucesor (3620, Nº 7806, h.\ 11); en mayo Julio González preside, como «\textit{Intendente Municipal}», la cooperadora asistencial del pueblo, con Mogro entre los vocales (4474, Nº 7843, h.\ 12), y el 4 de diciembre el gobernador declara intervenida la Municipalidad porque en ella «\textit{se habrían realizado patentamientos de automotores en forma irregular}», suspende a González, al que el decreto llama presidente de la comisión municipal, designa interventor a un suboficial principal retirado y manda instruir un sumario administrativo (7179, Nº 7961, h.\ 7). \textbf{En 1968}, con un gobernador designado en abril por el Poder Ejecutivo Nacional ---y con él la numeración de los decretos vuelve a empezar---, renuncia el interventor y se designa a Manuel David Serrey, que en agosto pasa a presidir también la cooperadora asistencial en lugar de González (57 y 58, Nº 8053, h.\ 14 y 15; 1408, Nº 8124, h.\ 8); el nombre es el de un miembro de la comisión del Cristo de 1964, David Serrey, y el acto no dice si es el mismo. En marzo, por terna de la Municipalidad intervenida, se designan por dos años un juez de paz titular y un suplente para cargos que el decreto da por vacantes (8596, Nº 8036, h.\ 8). Lo hallado de los dos años no trae el resultado del sumario. \textbf{En 1969}, también con pendientes, el gobernador cambia en agosto sin que la numeración de los decretos vuelva a empezar; en septiembre se acepta la renuncia del juez de paz titular designado en 1968, y lo hallado del año no trae quién lo reemplaza (6659, Nº 8399, h.\ 8), y en octubre el decreto que renueva la cooperadora asistencial llama a Serrey «\textit{Intendente Municipal}» (6882, Nº 8421, h.\ 12); si la intervención terminó entonces, lo hallado del año no lo dice. \textbf{En 1970}, también con pendientes, la Provincia le acepta el 30 de abril la renuncia como interventor y designa presidente de la comisión municipal a Justo Pastor Lizondo, que ese mismo día deja el juzgado de paz (8641, 8642 y 8643, Nº 8546, h.\ 12); en 1971 le rechaza a Lizondo la renuncia (335, Nº 8815, h.\ 9) y el decreto que renueva la cooperadora lo llama «\textit{Intendente Municipal de La Caldera}» (1797, Nº 8743, h.\ 5), y en noviembre de 1972 la Ley 4525 convoca para marzo de 1973 a elegir cuatro miembros de la comisión municipal de La Caldera y dos de la de Vaqueros, el municipio creado en 1970 (Nº 9154, h.\ 12; Ley 4527, Nº 9161, h.\ 21 y 22; capítulo \ref{cap:cot}). \textbf{En 1973}, el año de esa convocatoria, al presidente de la comisión lo sigue designando la Provincia: el 5 de junio le rechaza a Lizondo la renuncia, el 10 de julio se la acepta y el 13 de agosto designa presidente a Eusebio Correa, al que en julio de 1974 designa otra vez «\textit{por un nuevo período Constitucional}» (103, Nº 9289, h.\ 13; 529, Nº 9317, h.\ 10; 819, Nº 9335, h.\ 10; 4981, Nº 9571, h.\ 8 y 9); en Vaqueros, el 4 de junio acepta la renuncia de Catalán Arellano, con las de otros treinta y un presidentes de comisión, y el 6 designa a Werfil Gallo ---el nombre del que presidió la comisión municipal de La Caldera de 1958 a 1962; el acto no dice si es el mismo--- (80, Nº 9289, h.\ 6 y 7; 129, Nº 9292, h.\ 10). Ningún acto hallado de 1973 da el resultado de la elección convocada para el 11 de marzo. Ese año, por terna de la Municipalidad, se designa juez de paz titular a Teodoro Bartolo Mogro y suplente a Liborio Gerardo Guerrero «\textit{por un nuevo período legal de dos años}»; en julio se le acepta la renuncia a Mogro, y en enero de 1974 el titular designado es Gerardo Portal (121, Nº 9218, h.\ 10 y 11; 736, Nº 9327, h.\ 13; 3237, Nº 9544, h.\ 5). Y el auxiliar de enfermería del puesto sanitario de Vaqueros tiene licencia desde mayo de 1973 «\textit{mientras se desempéñe como Legislador (Diputado Provincial)}» \textit{[sic]} (1265, Nº 9363, h.\ 9). \textbf{En 1975}, también con pendientes y con la provincia bajo intervención federal desde fines de 1974, la Provincia le acepta a Correa la renuncia el 8 de enero, en el mismo tramo de decretos en que acepta la del presidente de la comisión municipal de El Bordo, y el 31 designa a Luis Xamena «\textit{Interventor de la Municipalidad de La Caldera}», que en septiembre firma como «\textit{Interventor Municipal}» la única ordenanza de La Caldera que aparece en lo hallado del año (2, Nº 9667, h.\ 5; 92, Nº 9692, h.\ 7; 3519, Nº 9877, h.\ 12; capítulo \ref{cap:infraestructura}, más adelante); en Vaqueros, en julio, acepta la renuncia de Werfil Gallo y designa presidente de la comisión municipal a José Santiago Catalán, con el apellido del primer presidente de esa comisión (1764, Nº 9802, h.\ 6), y acepta también la renuncia de Serapio Guantay, «\textit{Juez Titular de Vaqueros}», sin que lo hallado del año le nombre sucesor (1383, Nº 9784, h.\ 6). \textbf{En 1976}, también con pendientes, la provincia pasa el 24 de marzo de la intervención federal a un «\textit{Interventor Militar}» y en abril a un gobernador designado, y los dos municipios del departamento quedan a cargo de suboficiales retirados: el 31 de marzo la Provincia acepta la renuncia de Xamena y designa interinamente en su lugar al suboficial mayor (R) Armando Fernández; el 3 de mayo dispone el cese de José Santiago Catalán en Vaqueros y designa interinamente al suboficial mayor (R) Telmo Camacho, que tendrá a su cargo, además de las tareas del ejecutivo, «\textit{aquellas inherentes al Concejo Deliberante}» ---la fórmula de los decretos del mismo día para otros municipios---, y el 23 de julio confirma a los dos, con el «\textit{personal militar que seguidamente se detalla}» al frente de otros nueve, como presidentes de sus comisiones municipales (116, Nº 9967, h.\ 10; 513, Nº 9982, h.\ 6; 1586, Nº 10049, h.\ 8 y 9). En noviembre renuncia el juez de paz suplente de La Caldera, Liborio Gerardo Guerrero, con el nombre del suplente designado en 1973, y lo hallado del año no trae a su sucesor (3016, Nº 10121, h.\ 15). \textbf{En 1977 y 1978}, también con pendientes, lo hallado no trae ninguna designación nueva al frente de los dos municipios: Fernández firma como presidente de la comisión municipal las ordenanzas de La Caldera de febrero de 1977 y de marzo de 1978, con dos secretarios generales distintos ---el segundo, otro suboficial mayor retirado---, y Camacho las de Vaqueros (2172, Nº 10282, h.\ 5 y 6; 2848, Nº 10316, h.\ 28 y 29; 1121, Nº 10533, h.\ 8; 3551, Nº 10375, h.\ 10 y 11; 1470, Nº 10577, h.\ 7), y el 19 de abril de 1977 el gobernador Gadea entrega el mando a otro capitán de navío retirado, Roberto Augusto Ulloa (1270, Nº 10230, h.\ 9). En julio de 1977 la Provincia designa juez de paz titular de La Caldera, a propuesta de la comuna, a Hugo Justino Muñoz, al que en febrero de 1978 le acepta la renuncia sin que lo hallado del año le nombre sucesor, y en septiembre de 1977, por terna de la Municipalidad de Vaqueros, a los jueces de paz titular y suplente de esa localidad, Salvador del Valle Mangogna ---con el nombre, escrito Mangogña, del juez de paz que en 1965 sale de la terna de un concejo deliberante; los actos no dicen si es el mismo--- y Benito Antonio Mangogna (2426, Nº 10294, h.\ 8 y 9; 236, Nº 10436, h.\ 10; 3081, Nº 10334, h.\ 22). \textbf{En 1979}, también con pendientes, un decreto de febrero ratifica las licencias que la Secretaría de Estado de Municipalidades había dado desde el 27 de noviembre de 1978 a Fernández y a Camacho, convocados al servicio activo por «\textit{el reciente operativo de movilización efectuado por las Fuerzas Armadas del País}», y el encargo del despacho a quienes los seguían en jerarquía: en La Caldera, el secretario general, otro suboficial mayor retirado, y en Vaqueros, el secretario-tesorero; las resoluciones que ratifica llaman a los dos intendentes municipales (195, Nº 10702, h.\ 5). Cuándo volvieron no lo dice ningún acto hallado. El 28 de noviembre de 1979 la Provincia le acepta a Fernández la renuncia a la presidencia de la comisión municipal, por «\textit{razones de índole personal}», y designa en su lugar a Diego Jorge Montero, que ejercerá, además de las funciones del ejecutivo, «\textit{todas aquellas facultades inherentes al Concejo Deliberante}» (1675 y 1676, Nº 10875, h.\ 7 y 8). El juzgado de paz de Vaqueros pierde a sus dos jueces entre marzo y abril, recibe en junio, por dos años, a Digna del Carmen Eufemia Robertson de Martínez y a Juan de la Cruz Díaz, y los vuelve a perder en agosto y en noviembre sin que lo hallado del año les nombre sucesores; el de La Caldera recibe en octubre a José Félix Robles como titular y a un Esteban Mogro como suplente, sin que el decreto diga a quiénes reemplazan (304, Nº 10714, h.\ 15; 354, Nº 10718, h.\ 7 y 8; 672, Nº 10756, h.\ 27; 1126, Nº 10821, h.\ 16; 1572, Nº 10866, h.\ 7; 1450, Nº 10849, h.\ 10). \textbf{En 1980}, también con pendientes, la Provincia cubre en marzo, a propuesta de la Municipalidad, los dos cargos de juez de paz de Vaqueros, vacantes, con Rubén Amado Fiori como titular y Neit Akrit Ferro Podestá como suplente (434, Nº 10953, h.\ 7), y el 1 de diciembre le acepta a Camacho la renuncia a la presidencia de la comisión municipal de Vaqueros y designa a otro suboficial mayor retirado, Julio Francisco Michaux, con las facultades del Concejo Deliberante, como a Montero el año anterior (1712 y 1713, Nº 11126, h.\ 27 y 7). En La Caldera, Montero firma ese año las ordenanzas que dan a una calle el nombre de Almirante Guillermo Brown y crean el escudo del municipio, con un secretario administrativo nuevo, Norberto Eduardo Camaño; y con su firma se publica también el presupuesto de 1979, fechado el 27 de julio de 1979, cuatro meses antes de su designación, sin que el decreto que lo aprueba lo explique (1687, Nº 11125, h.\ 9 y 10; 1769, Nº 11133, h.\ 4 a 6; 1146, Nº 11059, h.\ 6 y 7). El escudo lo elige el 30 de junio un «\textit{Consejo Asesor Comunal}» que Montero convoca en su despacho, integrado por Arturo Renée Fernández, Esteban Mogro, José Félix Robles y Esteban Muñoz ---Robles y Mogro con los nombres de los jueces de paz designados en 1979, y Fernández con el del presidente del Centro Gaucho que se sigue enseguida, escrito René; que sean las mismas personas no consta--- (1769, Nº 11133, h.\ 5). \textbf{De 1981 a 1983}, leídos con la misma reserva, los dos municipios cambian tres veces de presidente, y las tres por decreto. En Vaqueros, la Provincia le acepta a Michaux la renuncia en abril de 1981 y designa interinamente a un ingeniero de la Secretaría de Estado de Municipalidades, Sergio Horacio Cristóbal Oller, también con las facultades del Concejo Deliberante; en junio le da por cumplida la gestión, para devolverlo a su repartición, y designa titular a Enrique Gustavo Clement (388 y 389, Nº 11212, h.\ 7 y 8; 867 y 868, Nº 11263, h.\ 7): cuatro presidentes en doce meses, contando a Camacho. En La Caldera, la Provincia le acepta a Montero la renuncia desde el 1 de septiembre de 1982, y ese día el decreto 877 designa presidente a Arturo Renée Fernández ---René en los decretos de 1983---, desde que tome posesión, con las funciones del ejecutivo y las del Concejo Deliberante (876 y 877, Nº 11555, h.\ 5 y 6): es el decreto que invocan las ordenanzas de 1983. En marzo de 1983 la Provincia rechaza las renuncias de los jefes comunales de cuarenta y cinco municipios, entre ellos Fernández y Clement, y los confirma (360, Nº 11692, h.\ 8 y 9), y el 9 de diciembre les acepta las de los dos desde el 11 (2050 y 2056, Nº 11870, anexo, h.\ 50). El juzgado de paz de La Caldera pierde a su suplente en mayo de 1981, cuando la Provincia le acepta la renuncia a Mogro (669, Nº 11236, h.\ 8); en noviembre de 1982, a propuesta en terna de la Municipalidad y por vencimiento del período legal, designa por dos años a Mogro como titular y a Robles como suplente (1258, Nº 11609, h.\ 13 y 14), y en noviembre y diciembre de 1983 les acepta la renuncia a los dos (1872, Nº 11868, h.\ 8; 2013, Nº 11870, anexo, h.\ 45). Y dos nombres de estos años reaparecen en la ordenanza de 1983 que se transcribe más abajo: Esteban Mogro, el juez de paz suplente que la Municipalidad propone en terna en 1961 y que renuncia en febrero de 1962 (20.209-G, Nº 6490, h.\ 11; 1204-G, Nº 6567, h.\ 13), y al que en julio de 1963, por terna de la Intervención Municipal, se designa juez de paz propietario por dos años (8410, Nº 6908, h.\ 11), que renuncia desde enero de 1964, vuelve al juzgado en junio y lo deja en mayo de 1965 por la presidencia de la comisión municipal (1540, Nº 7020, h.\ 5), es el de su secretario administrativo, que en noviembre de 1983 renuncia además como juez de paz titular, y Arturo René Fernández, a quien la Provincia presta en 1962 un toro para la finca Campo Alegre (3634-E, Nº 6671, h.\ 12) y que en 1965 preside el Centro Gaucho «El Palenque» del pueblo (8792, Nº 7357, h.\ 5) ---en 1969 un Arturo René Fernández, agricultor, integra el órgano de fiscalización de la cooperadora asistencial, en la que es vocal un Esteban Mogro, comerciante (6882, Nº 8421, h.\ 12), y un Arturo Fernández firma como presidente la convocatoria de El Palenque (aviso 2068, Nº 8455, h.\ 40), y en 1973 firma así una convocatoria y con el nombre completo, Arturo René Fernández, la siguiente (avisos 14792, Nº 9287, h.\ 35; 14956, Nº 9298, h.\ 16)---, coincide con el de su presidente, que firma Arturo R.\ Fernández. Que sean las mismas personas no consta. \textbf{En 1984 y 1985}, leídos con la misma reserva, la designación sigue: el 16 de enero de 1984 el gobierno constitucional designa, desde que tomen posesión e invocando el artículo 178 de la Constitución, a Mario Hipólito García presidente de la comisión municipal de La Caldera y a Jorge Antonio de Bustos presidente de la de Vaqueros (145 y 146, Nº 11899, h.\ 6 y 7), y un decreto del 28 de diciembre, publicado en enero de 1985, los vuelve a designar, con los de otros veintitrés municipios, «\textit{por el período que fija el artículo 182º de la Constitución Provincial y artículo 35 de la Ley nº 1349}» (2862, Nº 12134, h.\ 5). En julio de 1985 la convocatoria para el 3 de noviembre vuelve a darles a los dos municipios dos miembros titulares y dos suplentes de sus comisiones, por dos años, y no el presidente (1465, Nº 12268, h.\ 5 y 6); la lista del Frente Justicialista lleva a De Bustos como candidato a convencional constituyente titular y a García como suplente, y el resultado no se publica en lo hallado del año (aviso 57995, Nº 12315, h.\ 28 y 29). Los jueces de paz se siguen designando por terna municipal, y tres de los cuatro designados en 1984 se van antes de terminar su período: en marzo de 1984 la Provincia designa por dos años a Juan García y a Enrique Humberto González en La Caldera y a Lía Hortensia Quipildor y a Hilda Ester Perales de Mangoña en Vaqueros (545, Nº 11935, h.\ 7; 576, Nº 11938, h.\ 4); en mayo le acepta la renuncia a la titular de Vaqueros, y en julio de 1985 a la suplente, y en agosto designa, por terna de la comisión municipal, a Rubén Amado Fiori y a Danton Cermesoni en los dos cargos, que el decreto da por vacantes (947, Nº 11972, h.\ 5; 1360, Nº 12264, h.\ 9; 1709, Nº 12292, h.\ 7 y 8); y el 30 de septiembre de 1985, a pedido de la Municipalidad, le acepta la renuncia al titular de La Caldera, sin que lo hallado del año le nombre sucesor (1966, Nº 12323, h.\ 5). Danton Julio Cermesoni es además el senador suplente de esa lista por el departamento, y Dantón Cermesoni, el nombre de un candidato a constituyente por Chicoana en 1948 (capítulo \ref{cap:siglo}); que sean la misma persona no consta. Y los dos nombres de siempre siguen apareciendo: en noviembre de 1984 un Arturo Fernández, sin segundo nombre, firma como presidente la convocatoria de la Agrupación Tradicionalista Gauchos de Güemes «El Palenque» (aviso 54818, Nº 12096, h.\ 17 y 18), y en febrero de 1985 un Esteban Mogro preside el comité departamental de un partido que sólo una anotación a mano del ejemplar escaneado llama «Unión Provincial» (aviso 55551, Nº 12156, h.\ 21 y 22). Que sean las mismas personas tampoco consta. \textbf{En 1986 y 1987, con la misma reserva, la serie de las designaciones se completa}. En febrero de 1986 el decreto 229 designa por un nuevo período a los presidentes de las comisiones municipales que lista su anexo, García en La Caldera y De Bustos en Vaqueros, los dos con el título del 2862 de 1984, y en enero la Provincia le acepta la renuncia al juez de paz suplente de La Caldera designado en 1984 (229, Nº 12398, h.\ 5; 128, Nº 12392, h.\ 5); lo hallado de 1986 a 1988 no le nombra sucesor ni a él ni al titular que renunció en 1985. El 13 de mayo de 1987 el decreto 907 le acepta a García la renuncia ---cita como su título sólo el 145 de 1984--- y el 908, del mismo día y del mismo expediente, invocando el artículo 164 y las cláusulas segunda y décima de las disposiciones transitorias de la Constitución, designa presidente de la comisión municipal a Teodoro Bartolo Mogro; que reemplace a García lo dicen la fecha y el expediente, no el decreto (Nº 12710, h.\ 8). Tres semanas después el decreto 1080 convoca para el 6 de septiembre a elegir, por cuatro años, intendentes en cincuenta y ocho municipios, La Caldera y Vaqueros entre ellos, y, por dos, tres concejales titulares y tres suplentes en cada uno de los dos (Nº 12721, h.\ 4 y 5): \textbf{el titular del ejecutivo municipal vuelve a estar en la boleta}. Quién resultó electo no lo publica lo hallado de 1987 ni el de 1988; en febrero de 1988 el decreto 271 da por cesados desde el 11 de diciembre de 1987 a los intendentes de treinta y ocho municipios que no fueron reelectos, y ni La Caldera ni Vaqueros están en la lista ---tampoco Cachi, Metán ni La Poma---, de modo que el decreto no permite saber quién siguió (Nº 12920, h.\ 5). Los partidos se habían adelantado: en octubre de 1986 la Unión Cívica Radical llama a internas para elegir candidatos a intendente en La Caldera y en Vaqueros, municipios que el decreto 229 trataba como comisiones con presidente, y el convenio de defensas que la Municipalidad de Vaqueros firma en junio de ese año ya llama intendente a De Bustos (Res.\ 9/86, Nº 12567, h.\ 19 a 21; 2447, Nº 12551, h.\ 6 y 7). Mogro es también el apellido del secretario administrativo de la comisión de 1983 (más abajo); que haya parentesco no consta.

\subsection*{No se elegía el ejecutivo municipal}

\begin{ficha}{Ley 6131, promulgada el 30 de julio de 1983 --- B.O.\ Nº 11.780, h.\ 4 y 5}
Convoca al pueblo de la Provincia a los comicios generales del \textbf{30 de octubre de 1983}.

\textbf{Art.\ 3º:} los departamentos de Cachi, Cafayate, Guachipas, Iruya, \textbf{La Caldera}, La Candelaria, La Poma, La Viña, Los Andes, Molinos, San Carlos y Santa Victoria eligen \textbf{un diputado titular y un suplente, y un senador titular y un suplente}. Es la representación mínima de la escala.

\textbf{Art.\ 4º:} para determinar el número de diputados ``rige el último censo nacional practicado en el año 1980''.

\textbf{Art.\ 5º:} convoca al pueblo de los municipios ``a fin de elegir \textbf{Concejales Municipales y Miembros de Comisiones Municipales}''. Enumera cuatro grupos: once municipios que eligen nueve concejales --- Capital, Tartagal, Orán, Metán, entre otros ---, dieciocho que eligen cinco, y otros dos, de dieciséis y de trece municipios. \textbf{Ni La Caldera ni Vaqueros figuran en los dos primeros}: La Caldera está en el tercero, con cuatro miembros titulares y tres suplentes de comisión municipal, y Vaqueros en el cuarto, con dos y dos.
\end{ficha}

La convocatoria de 1983 no menciona intendentes. Convoca a elegir cuerpos deliberativos: concejales donde hay Concejo, miembros donde hay Comisión Municipal. El titular del ejecutivo municipal no estaba en la boleta.

Eso confirma lo que se desprende del artículo 22 de la Ley 1349 ---ver la advertencia del apéndice \ref{cap:normativa}--- y deja sin objeto la búsqueda de ``actas de proclamación de intendentes de 1983'': no existen porque no hubo elección de intendentes. \textbf{Y desplaza la pregunta}: lo que hay que buscar para ese período no son proclamaciones sino designaciones.

En septiembre la Ley 6163 corrige el número de La Caldera ---elegirá dos titulares y dos suplentes, «\textit{y no como se consignara en el inciso c)}»--- sin dar la razón (Nº 11811, h.\ 5 y 6). Los partidos, en cambio, se organizan por departamento: ese año la Unión Cívica Radical, el Partido Renovador, la Democracia Cristiana y el Partido Justicialista publican en el Boletín las autoridades de sus comités, juntas o consejos de La Caldera, con domicilios en La Caldera, La Calderilla, Vaqueros y la finca Lesser en la nómina justicialista (aviso 49513, Nº 11730, h.\ 43; aviso 50106, Nº 11768, h.\ 20 a 22; aviso 50292, Nº 11786, h.\ 13; aviso 50630, Nº 11806, h.\ 17; aviso 50656, Nº 11808, h.\ 17); presiden los comités radical y renovador y el consejo justicialista Salvador del V.\ Mangogña, Armando Fernández ---el nombre del presidente de la comisión municipal hasta 1979, sin que conste que sea el mismo--- y Martín Borja. Y un decreto reparte las quince mesas del departamento entre los circuitos 8 y 9, en escuelas de La Caldera, Potrero de Castilla, Los Yacones, Campo Alegre, Lesser y Vaqueros y en la estación de Mojotoro (1478, Nº 11826, h.\ 9 y 23).

\subsection*{Era una Comisión Municipal, y tiene nombre}

El Boletín Nº 11.870 transcribe íntegra una ordenanza de La Caldera y el decreto provincial que la aprueba.

\begin{ficha}{Ordenanza Nº 047 del 30 de junio de 1983}
Establece para ese año una reducción del cuarenta por ciento en las tasas de Alumbrado y Limpieza a favor de jubilados y pensionados, con tres condiciones: ser titular de la \textbf{única} propiedad, que el inmueble sea casa habitación del beneficiario, y que el haber jubilatorio no supere el mínimo de ley.

Firman: \textbf{Arturo R.\ Fernández\index{Fernández, Arturo R.}, Presidente de la Comisión Municipal de La Caldera}, y \textbf{Esteban Mogro\index{Mogro, Esteban}, Secretario Administrativo}.

Los considerandos invocan el \textbf{Decreto provincial 877/82}, que acuerda al titular del Departamento Ejecutivo Municipal, ``a más de las funciones que le son propias, \textbf{las inherentes al Concejo Deliberante}''.

La ordenanza dispone en su artículo 4º: ``Pase a Secretaría de Estado de Municipalidades \textbf{para su aprobación}''. Y en efecto, el \textbf{Decreto 1978} del 5 de diciembre de 1983, del Ministerio de Gobierno, Justicia y Educación, la aprueba y transcribe, cinco meses después de dictada (h.\ 22).
\end{ficha}

\begin{cautela}
Tres capas de tutela en un solo acto de gobierno local, y las tres importan.

\textbf{La primera: no había Concejo.} El Decreto 877/82 concentró en el titular del ejecutivo las funciones del cuerpo deliberativo. Quien sancionaba la ordenanza y quien la promulgaba eran la misma persona.

\textbf{La segunda: esa persona no era un intendente sino un Presidente de Comisión Municipal}, la figura que la Ley 1349 --- entonces vigente, derogada en 2018 --- reservaba en su artículo 33 a los municipios de tercera categoría, y cuyo titular ``será nombrado y removido por el Poder Ejecutivo''.

\textbf{La tercera: la ordenanza no valía por sí sola.} Necesitaba aprobación de la Secretaría de Estado de Municipalidades, y la obtuvo por decreto ministerial. Una rebaja de tasas para jubilados de un pueblo de mil habitantes recorrió el circuito completo hasta el Boletín Oficial de la Provincia.

\textbf{Y en 1921 la elección del departamento se anula entera.} La Junta Electoral resolvió, el 6 de septiembre, \textbf{declarar nula la elección del 7 de agosto en el distrito de la Caldera} y en otros once, «\textit{por no haberse verificado elección legal en dos tercios de las mesas}»; el decreto que convoca de nuevo es del 20 de septiembre y \textbf{se publica el 28 de octubre, diecinueve días después del segundo domingo de octubre que él mismo señalaba} (capítulo \ref{cap:siglo}). \textbf{Son dos años seguidos}: en 1920 se anuló la urna de Mojotoro y hubo complementarias; en 1921, el distrito completo. Ni la nota de la Junta ni la resolución de nulidad se publican, de modo que este libro no puede decir si la elección de octubre se hizo ni qué pasó con los dos convencionales que le correspondían al departamento en la reforma constitucional.

\textbf{Y en 1920 el instrumento que fija la representación no cierra consigo mismo.} El decreto 685 del 6 de febrero de 1920, que convoca a elegir legisladores provinciales, \textbf{nombra dos veces a Caldera con un diputado cada vez} en la misma enumeración, y enumera doce departamentos distintos, en trece renglones, de los veintidós (capítulo \ref{cap:siglo}). Ese mismo año, tres partidos proclaman candidato a diputado por el departamento, y uno de los tres ---\textbf{Dermidio Arancibia}--- es quien dos años antes había pedido las dos concesiones mineras que el Boletín registra sobre este suelo. \textbf{Coincidencia de nombre y no identificación registral}, pero es el primer caso de esta serie en que quien pide una concesión sobre el departamento se presenta a representarlo. Y la única urna anulada del departamento ese año es la de Mojotoro, que es donde está su única plaza rentada de policía, que cambió de manos tres veces en cinco meses.

\textbf{Y la tutela no empieza en 1982.} El capítulo \ref{cap:siglo} documenta que en marzo de 1918 el departamento fue incluido, por el decreto 1.557, entre los catorce que pasaban de Comisión Municipal a \textbf{Municipalidad electiva} por tener renta superior a cinco mil pesos; y que siete semanas después el decreto 1.660 lo sacó de la lista, y sólo a él, porque su cálculo de recursos había computado como renta un saldo del ejercicio anterior. \textbf{La renta real eran \$3.387,50}, y el decreto concluye que el departamento «\textit{no está en condiciones de ser administrado por una Municipalidad electiva}». La tercera categoría de la Ley 1349 es de 1933; el criterio de renta que la funda ya operaba en 1918, y ya dejaba a este departamento del lado de abajo. \textbf{Y diez meses después el gobierno siguiente cambió el fundamento sin cambiar el resultado}: el decreto 33 del 15 de enero de 1919, que nombra la Comisión Municipal del departamento, considera que el artículo de la Ley Orgánica en que se creaban municipalidades «\textit{ha sido derogado por el art.\ 175 de la Constitución, que exije para la existencia de Municipalidades autónomas un centro urbano de cinco mil habitantes, de que carecen todos los Departamentos de Campaña}». \textbf{El umbral deja de ser de pesos y pasa a ser de personas, y la cifra es la misma: cinco mil.} Es la que la Ley 1349 recogerá catorce años después como techo de la tercera categoría, y la que el capítulo \ref{cap:poblacion} encuentra todavía operando en la clasificación estadística de 2022.

De modo que el punto de partida del período que este libro estudia no es un municipio débil: es un municipio \textbf{tutelado}, sin órgano deliberativo, con ejecutivo designado y con sus actos sujetos a aprobación provincial.
\end{cautela}

\subsection*{Y confirma lo que el capítulo \ref{cap:poblacion} sospechaba}

El capítulo \ref{cap:poblacion} estableció que la clasificación M3 de la Dirección General de Estadísticas se presenta como ``municipio de tercera categoría'', que la tercera categoría de la Ley 1349 correspondía a las Comisiones Municipales, y que desde 2018 esa ley está derogada y el régimen vigente no tiene categorías. La categoría estadística, entonces, está \textbf{congelada}: sobrevive a la ley que la sostenía.

Estos documentos fechan el congelamiento. \textbf{En 1983 la clasificación era exacta}: ambos eran Comisiones Municipales. El Boletín Nº 11.870 registra que el 9 de diciembre de 1983 --- víspera del traspaso democrático --- el Poder Ejecutivo provincial aceptó las renuncias de los \textbf{Presidentes de las Comisiones Municipales} de una larga lista que incluye a \textbf{La Caldera y a Vaqueros}, y en decretos separados las de los \textbf{Intendentes de las Municipalidades}. Las dos figuras coexistían y el departamento estaba, entero, del lado de las Comisiones.

Con lo cual la categoría que hoy comparten con San Lorenzo --- treinta y dos mil habitantes --- no es un error de clasificación: es una clasificación de 1931 que en 1983 todavía era exacta, que nadie actualizó mientras la ley la exigía y que la estadística siguió usando después de derogada. Vaqueros tiene hoy 7.774 habitantes y sigue clasificado como lo estaba cuando tenía menos de dos mil: la Ley 4365 lo creó en 1970 como «municipio de tercera categoría (2da.\ clase)» (capítulo \ref{cap:cot}).

\begin{cautela}
Un dato menor que conviene registrar sin conclusión. El Boletín Nº 11.868 consigna que en noviembre de 1983 el Poder Ejecutivo aceptó la renuncia del señor \textbf{Esteban Mogro\index{Mogro, Esteban}} al cargo de \textbf{Juez de Paz Titular de la localidad de La Caldera}. Es el mismo nombre que cinco meses antes firmaba la Ordenanza 047 como Secretario Administrativo de la Comisión Municipal.

\textbf{Sucesivos no fueron}: el juez al que se le acepta la renuncia es el titular designado en noviembre de 1982 ---el decreto de 1983 cita el de 1982 y da la misma matrícula--- (1258, Nº 11609, h.\ 13 y 14; 1872, Nº 11868, h.\ 8), de modo que en junio y julio de 1983, cuando el secretario administrativo firma las ordenanzas 047 y 048, el juzgado de paz tenía un titular con su nombre. Si son la misma persona o dos homónimos en una localidad de mil habitantes, este libro no lo establece: las ordenanzas no dan su documento. Lo consigna porque, si es la misma persona, describe algo característico de las localidades pequeñas: la administración municipal y la justicia de paz atendidas por la misma mano. Con el mismo nombre hay un antecedente: en mayo de 1965 a un Esteban Mogro se le acepta la renuncia al juzgado de paz titular y a la semana la Provincia lo designa presidente de la comisión municipal (8430, B.O.\ Nº 7343, h.\ 9; 8560, Nº 7347, h.\ 18); ahí los cargos fueron sucesivos, y la presidencia termina en septiembre de 1966 (777, Nº 7666, h.\ 7). En febrero de 1967 un Esteban Mogro es designado juez de paz suplente, y en diciembre se le acepta la renuncia (3069, Nº 7774, h.\ 12; 7232, Nº 7964, h.\ 13). Y en octubre de 1979 un Esteban Mogro es designado otra vez juez de paz suplente de La Caldera, por dos años (1450, Nº 10849, h.\ 10), y en mayo de 1981 se le acepta la renuncia a ese juzgado suplente (669, Nº 11236, h.\ 8). En junio de 1980 un Esteban Mogro integra el Consejo Asesor Comunal que elige el escudo del municipio (1769, Nº 11133, h.\ 5).
\end{cautela}


\section{Lo que la banca dejó atrás}

El traslado de Moreno al Senado tuvo una consecuencia que este libro debe registrar, porque no ocurre en la Legislatura sino en el municipio donde se concentra la conversión de suelo del departamento.

Moreno renunció a la intendencia de Vaqueros faltando dos años de mandato. La regla aplicable no es la vieja Ley Orgánica de Municipalidades, derogada en 2018, sino el Régimen de Municipalidades de la \textbf{Ley 8126} (B.O.\ 17/12/2018). Su artículo 64 pone el Ejecutivo, en caso de impedimento, a cargo del presidente del Concejo Deliberante, cargo que el propio cuerpo elige cada año (art.\ 50); y su artículo 67 dispone que, si la renuncia es definitiva y falta \textbf{más de un año} para completar el mandato, \textbf{debe convocarse a elecciones} conforme el artículo 182 de la Constitución provincial. Faltaban dos. \pendiente{convocatoria a elecciones para completar el mandato de intendente de Vaqueros conforme el art.\ 67 de la Ley 8126, o acto que la descarte}

\begin{ficha}{Tres titulares del Ejecutivo en tres meses}
\begin{itemize}[leftmargin=*,nosep]
\item \textbf{Diciembre de 2025.} Moreno renuncia. Asume el Ejecutivo \textbf{Daniel Viveros}, concejal y entonces presidente del Concejo Deliberante.
\item \textbf{1 de marzo de 2026.} En la primera sesión del año, el Concejo elige presidenta a \textbf{Andrea Salvatierra}, lo que conforme a la normativa vigente en el municipio la habilita a asumir la intendencia interina. Viveros sostiene que no corresponde desplazarlo, porque la renuncia de Moreno fue definitiva y su asunción debe extenderse hasta una nueva convocatoria a elecciones.
\item \textbf{Primeros días de marzo de 2026.} Viveros recurre a la Justicia, denuncia un ``golpe institucional'' y solicita medida cautelar de no innovar. La \textbf{Corte de Justicia de Salta} rechaza el planteo: la elección de autoridades del Concejo es facultad legítima del cuerpo y no se acreditó irregularidad en el procedimiento. Salvatierra, que había jurado el 2 de marzo, queda confirmada; el fallo es del 3 de marzo.
\item \textbf{Mayo de 2026.} El senador Moreno ratifica públicamente su respaldo a la gestión de Salvatierra y marca distancia de la etapa anterior.
\end{itemize}
\textit{Fuente: cobertura de prensa provincial, diciembre de 2025 a mayo de 2026. Los textos de la resolución del Concejo y del fallo de la Corte no fueron consultados.}
\end{ficha}

\begin{cautela}
Lo que este episodio agrega no es color político. Es un dato de capacidad estatal, y es del mismo orden que los que el capítulo \ref{cap:redes} reúne.

Entre diciembre de 2025 y marzo de 2026, el municipio más presionado del departamento --- el que aprueba fraccionamientos, otorga factibilidades, emite certificados de aptitud ambiental y ejerce el poder de policía sobre la obra privada --- tuvo \textbf{tres titulares del Ejecutivo en tres meses}, y el tercero llegó tras una disputa que debió resolver la Corte provincial. La Ley 8126 sólo prevé la intervención de la Corte para el juicio político al intendente (art.\ 92); la vía por la que llegó este conflicto es la que la prensa describe, una cautelar de no innovar, y el fallo no se consultó.

Ese es también el período de la última reforma del Código de Planeamiento Urbano Ambiental de Vaqueros: la Ordenanza 648/26, promulgada por el Decreto 486/26 del \textbf{26 de enero de 2026}, que firma «J.\ Daniel Viveros, Intendente». La reforma que incorporó los distritos de la costanera se promulgó, entonces, durante el interinato que el Concejo daría por terminado cinco semanas después. El capítulo \ref{cap:vaqueros} analiza el texto que resultó de ella.

Corresponde una advertencia expresa, porque el material se presta a la insinuación y este libro no la hace: \textbf{no se encontró vinculación documental alguna entre la disputa por la presidencia del Concejo y materia urbanística}. Lo que se registra es la discontinuidad del Ejecutivo municipal en un trimestre determinado, no una hipótesis sobre su causa.
\end{cautela}


\section{El voto que impide la tesis fácil}

En diciembre de 2024, la actualización del Ordenamiento Territorial de Bosques Nativos fue aprobada por el Senado provincial con amplia mayoría y tres votos en contra: los senadores por Cachi, por San Martín y \textbf{por La Caldera}. Esa norma es la \textbf{Ley provincial Nº 8483}, publicada en el Boletín Oficial el 3 de enero de 2025, que revisa el ordenamiento, modifica las superficies por categoría de conservación e introduce la figura del Área de Producción y Conservación. Su contenido se analiza en el capítulo \ref{cap:opacidad}, más adelante. El senador por La Caldera es, además, el impulsor de la prórroga de la ley provincial de suspensión de desalojos de familias rurales hasta diciembre de 2026, norma que \textit{El dispositivo salteño} consignó en su apéndice sobre litigio territorial.

El registro nominal de esa votación es el contrapeso más fuerte que este libro tiene contra su propia tesis, y por eso va en el cuerpo del capítulo y no escondido en una nota. La representación departamental en la cámara que aprueba el mapa de bosques votó contra ese mapa.

De aquí se sigue una redirección del objeto. Si la banca senatorial no es el punto de captura, entonces la decisión efectiva sobre el uso del suelo ocurre donde este libro la encontró: la ordenanza del Concejo Deliberante, el certificado de aptitud ambiental municipal y el expediente de la Secretaría de Ambiente provincial. Tres instancias técnicas, de baja visibilidad pública, sin votación nominal publicada y sin cobertura periodística sistemática.

Ese desplazamiento --- del órgano político visible al órgano técnico opaco --- se desarrolla en el capítulo \ref{cap:opacidad}, más adelante, que es donde este libro reúne el conjunto de la evidencia sobre el punto.

\section{El proyecto de reserva, y la tensión que abre}

En marzo de 2023, siendo senador, Calabró presentó ante la Cámara un proyecto para crear un área protegida sobre la Ruta Nacional 9, en el norte del departamento.

\begin{ficha}{Expte.\ Nº 90-31.780/2023 --- Reserva Natural Corredor Biológico Abra de Santa Laura}
Presentado el \textbf{30 de marzo de 2023} por el senador \textbf{Héctor Miguel Calabró\index{Calabró, Héctor Miguel}}, girado a la \textbf{Comisión de Minería, Recursos Naturales y Medio Ambiente}.

\textbf{Art.\ 1º:} instituye la Reserva Natural Corredor Biológico Abra de Santa Laura en el departamento La Caldera, ``de acuerdo al contenido de la \textbf{Ley Provincial 4.736}''.

\textbf{Art.\ 3º --- fijación de límites:} ``se efectuará sobre el eje del \textbf{talweg del río Porongal}, y su área de influencia de media y alta Sub Cuenca'', conforme los estudios y mensuras que establezca el Poder Ejecutivo por decreto, que además debe designar autoridad de aplicación entre los Ministerios de Producción, Ambiente y Turismo. Dentro de los noventa días de aprobado, debe elaborarse la reglamentación de funcionamiento.

\textbf{Art.\ 4º:} invita a asesorar a la Universidad Nacional de Salta --- Facultades de Ciencias Naturales y de Economía --- y a la Universidad Católica de Salta, Escuela de Turismo.

\textbf{Art.\ 5º --- Dominio Público:} ``Se declaran de Dominio Público Provincial las tierras comprendidas en este Corredor Biológico, conforme el régimen jurídico sustentado en \textbf{acuerdos con los propietarios} de las mismas, \textbf{sin necesidad de expropiaciones}'', por analogía con la Ley 6.806/95 de la Reserva Natural de la Quebrada de las Conchas.

\textbf{Destino: caducó.} Consta en el registro de proyectos caducados de la Cámara, por aplicación de la Ley 1111, con fecha \textbf{13 de marzo de 2025}. Nunca fue tratado.
\end{ficha}

\begin{cautela}
Cinco observaciones sobre el texto, en orden de importancia.

\textbf{Primera: caducó.} El proyecto entró en marzo de 2023, fue girado a comisión y ahí quedó. La Ley 1111 lo dio por caduco en marzo de 2025, nueve meses antes de que terminara el mandato de su autor. La reserva no existe y nunca estuvo cerca de existir.

\textbf{Segunda: no invoca la ley vigente.} El artículo 1º remite a la Ley provincial 4.736, no a la Ley 7107 del Sistema Provincial de Áreas Protegidas, sancionada en 2000 y vigente. Este libro \textbf{no afirma que la 4.736 esté derogada}: el texto completo de la Ley 7107, obtenido y leído, \textbf{no contiene artículo derogatorio alguno} --- su artículo 63 se limita a remitir supletoriamente a la Ley 7070, y su artículo 35 inciso c faculta a ``recategorizar si correspondiere, las áreas protegidas existentes en la Provincia al momento de la promulgación''. Si la 4.736 subsiste, en qué medida y con qué alcance, es una cuestión que excede a este trabajo.

Lo comprobable es que el proyecto se apoyó en una norma anterior al régimen integral vigente, y no en el régimen integral vigente --- que es, además, el que el capítulo \ref{cap:tierra} identificó como el instrumento disponible y no usado en el departamento.

\textbf{Tercera: la ley no era la única vía, ni la más rápida.} Los artículos 60 y 61 de la Ley 7107 permiten declarar un área protegida por \textbf{acto administrativo del Poder Ejecutivo}, de oficio o a petición de parte, previa evaluación de la Autoridad de Aplicación. Un proyecto de ley que caduca en comisión no agota el asunto: la vía administrativa quedaba abierta antes, durante y después. Y una vez declarada por decreto, el área sólo puede alterarse por ley --- artículo 62 ---, de modo que esa vía produce además un resultado más difícil de revertir.

\textbf{Cuarta: el artículo 5º del proyecto es jurídicamente singular.} Declara de dominio público provincial tierras privadas ``conforme acuerdos con los propietarios'' y ``sin necesidad de expropiaciones''. El dominio público no se constituye por acuerdo. El artículo 51 de la Ley 7107 prevé lo contrario: obliga a impulsar expropiación cuando las restricciones impidan el ejercicio regular de los derechos. Y su artículo 53 ofrece la figura que el proyecto parecía buscar sin nombrarla --- la incorporación voluntaria del propietario, con beneficios impositivos y renuncia no admisible antes de veinte años ---, que no requiere declarar público lo privado.

\textbf{Quinta: la geografía del texto no coincide con la de la prensa.} La cobertura de julio de 2023 ubicó la reserva entre los kilómetros 1.642 y 1.650 de la Ruta 9, del Dique Campo Alegre al límite con Jujuy. El articulado la ubica sobre el eje del \textbf{talweg del río Porongal} y su subcuenca media y alta. Pueden ser la misma zona descripta de dos maneras --- una por la ruta, otra por la cuenca ---, y este libro no afirma que no lo sean. Cabe consignar que el instrumento delimita por hidrografía y la difusión por kilometraje vial.
\end{cautela}

\begin{cautela}
Un nombre que conviene retener sin sacar conclusiones. La Ley 4486 --- analizada en el capítulo \ref{cap:expropiacion} --- expropió en 1972 una fracción de la finca \textbf{Los Porongos}, que el artículo 6º del PDUA sigue enumerando entre las unidades del municipio. El proyecto de reserva delimita por el \textbf{río Porongal}. La proximidad toponímica es evidente y puede no significar nada: los topónimos se repiten sobre cuencas y las fincas suelen llevar el nombre del curso que las atraviesa.

Establecer si hay relación --- y de haberla, cuál --- requiere el antecedente dominial de esa finca y la cartografía de la subcuenca. Es una de las cuatro o cinco preguntas que este libro deja abiertas sobre el mismo hueco: quién es dueño de qué en el departamento.
\end{cautela}

\begin{cautela}
Corresponde ahora enunciar, con la mayor precisión posible, una tensión que este libro documenta y no resuelve.

El capítulo \ref{cap:loteo} establece, con avisos del Boletín Oficial, que \textbf{Héctor Miguel Calabró\index{Calabró, Héctor Miguel} ---identificado por el mismo número de documento que en el acta de proclamación de 2011--- fue designado gerente suplente de Jardín Celestial S.R.L.} por acta del 8 de noviembre de 2018, publicada en enero de 2020, sin que este relevamiento haya encontrado constancia posterior de su cese ni de su continuidad; y que esa sociedad es copropietaria del inmueble sobre el que se aprobó el plano del loteo El Durazno, en el sur del pueblo de La Caldera. Este capítulo establece que, en marzo de 2023, el mismo Calabró impulsó como senador un área protegida para frenar el avance inmobiliario, en el norte del departamento.

\textbf{Las dos cosas no se contradicen en el plano de los hechos.} La reserva propuesta y el loteo están en extremos opuestos del departamento, separados por el pueblo y por el embalse; proteger un tramo de Yungas sobre la cornisa no impide nada en la zona sur. Un legislador puede impulsar conservación en un lugar y tener interés societario en otro sin incurrir en irregularidad alguna, y este libro no afirma que la haya.

Lo que efectivamente queda registrado, y es un dato estructural, es que \textbf{la misma persona fue designada en 2018 en el órgano de administración de una sociedad desarrolladora del departamento ---designación publicada en 2020, sin constancia posterior de cese--- y ocupó entre 2021 y 2025 la banca que legisla sobre el uso del suelo de ese departamento}, y que desde esa banca votó en diciembre de 2024 contra la Ley 8483 e impulsó en 2023 un área protegida.

Este libro no atribuye intención a ninguno de esos actos. Enuncia la superposición de posiciones porque es verificable en instrumentos públicos, y porque la superposición de posiciones --- no la mala fe --- es exactamente el objeto que ambos libros describen. La verificación que falta, y que es previa a cualquier publicación, es el derecho a réplica: si seguía siendo gerente suplente durante su mandato de senador y, en ese caso, qué dice el propio Calabró sobre la compatibilidad entre ambas posiciones, si informó la participación societaria en las declaraciones juradas patrimoniales que la normativa exige a los legisladores, y si se excusó en alguna votación relativa al departamento.
\end{cautela}


\section{La heterogeneidad del poder local}

El libro documentó, sobre el mismo cuerpo deliberativo y en la misma semana de enero de 2026, dos actos de signo opuesto: el pedido de clausura preventiva del ingreso a un loteo cerrado y la promulgación de una ordenanza, sancionada a fines de 2025, que prohíbe cobrar el acceso al perilago de un embalse público. Documentó también que el propio municipio denunció y logró la clausura de un loteo ilegal en 2020, y que su Concejo pidió la clausura de dos campings en 2026.

La descripción del dispositivo caldereño debe incluir su heterogeneidad o no describe nada. No hay aquí un poder local capturado enfrentado a una sociedad civil indefensa: hay un aparato que denuncia y habilita, que clausura y agradece la inversión, y cuyas fisuras internas son tan documentables como sus continuidades.

Lo que este libro sostiene no es que alguien haya capturado al Estado local. Es que se le asignó una función sin los medios para ejercerla, y que el resultado de esa asimetría se distribuye de manera muy desigual entre quienes habitan el territorio.

\begin{cautela}
Sobre la serie de intendencias, este relevamiento avanzó de manera desigual y conviene separar los tramos por su grado de verificación.

\textbf{Verificado en fuente oficial.} Se obtuvieron las actas de proclamación de electos del Tribunal Electoral de la Provincia correspondientes a las elecciones generales del \textbf{10 de abril de 2011} y del \textbf{17 de mayo de 2015}. Ambas consignan, para el departamento La Caldera:

\begin{center}\small
\begin{tabular}{@{}lll@{}}
\toprule
\textbf{Año} & \textbf{Cargo} & \textbf{Electo} \\
\midrule
2011 & Diputado & Héctor Miguel Calabró\index{Calabró, Héctor Miguel} (PJ) \\
2011 & Intendente de La Caldera & \textbf{Luis Gerardo Mendaña\index{Mendaña, Luis Gerardo}} (PJ) \\
2011 & Intendente de Vaqueros & Daniel Roberto Moreno Ovalle\index{Moreno Ovalle, Daniel Roberto} (SST) \\
2015 & Intendente de La Caldera & \textbf{Daniel Alejandro Escalera\index{Escalera, Daniel Alejandro}} (PJ) \\
2015 & Intendente de Vaqueros & Daniel Roberto Moreno Ovalle\index{Moreno Ovalle, Daniel Roberto} (PV) \\
\bottomrule
\end{tabular}
\end{center}

\textbf{La contradicción quedó resuelta.} Un listado de segunda mano atribuía la intendencia electa en 2011 a ``Luis Gerardo Escalera''. El acta oficial dice \textbf{Luis Gerardo Mendaña\index{Mendaña, Luis Gerardo}}. El error del listado consistía en unir el nombre de pila de uno con el apellido del otro --- Escalera es el electo en 2015, no en 2011 ---, y es un buen recordatorio de por qué este libro no publica lo que no verifica.

Del tramo anterior, \textbf{Héctor Miguel Calabró\index{Calabró, Héctor Miguel} fue intendente entre 1999 y 2011}, en tres mandatos consecutivos, sucediendo a \textbf{Horacio Martín Quipildor}. Eso no consta en actas obtenidas sino en tres fuentes periodísticas coincidentes, incluida una nota de Página/12 de 2021 donde el propio Calabró declara que en 2010 ``era intendente''. Varios decretos provinciales lo confirman para el tramo final ---el 1639/07, el 4354/09 y el 556/10 aprueban convenios que firma como intendente---, entre ellos dos que lo nombran expresamente en ese carácter: el Decreto 3318/08 (B.O.\ Nº 17.931, 19/08/2008) lo designa presidente de la Comisión Directiva de la Cooperadora Asistencial de La Caldera en su carácter de «Intendente Municipal», y el Decreto 2508/09 (B.O.\ Nº 18.129, 18/06/2009), que aprueba un convenio con el PAMI para la refacción y ampliación del Hogar de Ancianos del pueblo, lo nombra en el Comité Coordinador como «Intendente Municipal de la localidad de La Caldera». Y uno de 2002 lo pone en funciones en el primer mandato: el Decreto 1106 del 8 de julio de ese año, que traslada simbólicamente la capital de la Provincia a La Caldera para el 9 de julio, pone a cargo del «Intendente de La Caldera, Dr.\ Héctor Miguel Calabró» la entrega del decreto que declara huésped de honor al gobernador (B.O.\ Nº 16434, h.\ 7 y 8). El inicio del mandato en 1999 y las dos reelecciones siguen apoyados en la prensa. De su antecesor, el Boletín Oficial de 1990 registra otro cargo: en noviembre de ese año Horacio Martín Quipildor figura como secretario del Concejo Deliberante de la Municipalidad de La Caldera ---y con él un concejal del mismo municipio--- en el anexo del decreto 2367, que da de baja de sus cargos provinciales a los agentes con doble empleo (B.O.\ Nº 13558, h.\ 69 y 132).

\textbf{No corroborado, y por lo tanto no publicable.} Ese mismo listado propone para 1983, 1987 y 1991 tres nombres --- Roberto Adán Galli, Alberto Javier Alderete y Víctor Abelardo Montoya --- que ninguna fuente independiente respalda. \textbf{Quedan fuera de la serie} hasta que aparezcan en las proclamaciones del Tribunal Electoral, que es donde constan.
\end{cautela}

\begin{ficha}{Intendencias de La Caldera --- lo verificado}
\begin{center}\small
\begin{tabular}{@{}llp{4.2cm}@{}}
\toprule
\textbf{Período} & \textbf{Titular} & \textbf{Estado} \\
\midrule
1982--1983 & Arturo René Fernández & \textbf{Boletín Oficial}: presidente de la comisión municipal, designado por decreto y con renuncia aceptada por decreto \\
ene.\ 1984--may.\ 1987 & Mario Hipólito García & \textbf{Boletín Oficial}: presidente de la comisión municipal, designado por decreto (145 de 1984, renovado por el 2862 de 1984 y el 229 de 1986), con renuncia aceptada por decreto (907 de 1987) \\
may.\ 1987 & Teodoro Bartolo Mogro & \textbf{Boletín Oficial}: presidente de la comisión municipal, designado por decreto (908 de 1987) \\
1987--1999 & --- & Sin verificar: la convocatoria para el 6 de septiembre de 1987 incluye intendente por cuatro años, y el resultado no está en lo hallado de 1987 y 1988 \\
hasta 1999 & Horacio Martín Quipildor & Prensa coincidente \\
1999--2011 & Héctor Miguel Calabró\index{Calabró, Héctor Miguel} & Prensa coincidente (tres mandatos); \textbf{Boletín Oficial} para 2002 y para 2007 a 2010 \\
2011--2015 & Luis Gerardo Mendaña\index{Mendaña, Luis Gerardo} & \textbf{Acta de proclamación} \\
2015--2019 & Daniel Alejandro Escalera\index{Escalera, Daniel Alejandro} & \textbf{Acta de proclamación} \\
2019--2027 & Diego Ramiro Sumbay\index{Sumbay, Diego Ramiro} & Prensa coincidente; actas no obtenidas \\
\bottomrule
\end{tabular}
\end{center}
\end{ficha}

\begin{cautela}
La serie verificada obliga a registrar una coincidencia temporal, y este libro la enuncia con el mismo cuidado con que enunció las anteriores.

Jardín Celestial S.R.L.\ se constituyó por instrumento publicado en el Boletín Oficial del \textbf{1 de diciembre de 2010}. El capítulo \ref{cap:loteo} documenta que entre sus socios fundadores figura \textbf{Carlos Miguel Calabró\index{Calabró, Carlos Miguel}}, entonces de diecinueve años y estudiante, con el treinta por ciento del capital. En esa fecha, \textbf{Héctor Miguel Calabró\index{Calabró, Héctor Miguel} era intendente de La Caldera} y lo sería hasta diciembre de 2011.

Es decir: la sociedad que años después sería copropietaria del inmueble del loteo El Durazno se constituyó, con un socio de ese apellido, durante el tercer mandato del intendente del municipio donde la sociedad tenía su objeto.

\textbf{Este libro no sostiene que exista relación entre ambos hechos.}

Sobre el parentesco, este libro se atiene a lo que puede acreditar. El vínculo entre Héctor Miguel Calabró\index{Calabró, Héctor Miguel} y Carlos Miguel Calabró\index{Calabró, Carlos Miguel} es \textbf{altamente probable} --- apellido y segundo nombre compartidos, ambos domiciliados en La Caldera en 2010, diferencia de edad compatible --- pero \textbf{no está acreditado con partida del Registro Civil}, y por eso ni este capítulo ni el \ref{cap:loteo} lo afirman. Se acredita con una partida, no con un apellido, y es una de las preguntas del pedido de derecho a réplica. Tampoco afirma que la constitución de una sociedad de cementerio parque durante una intendencia tenga nada de irregular, porque no lo tiene.

Lo que corresponde registrar es la secuencia, porque es verificable y porque el libro registró todas las demás: 2010, constitución con un socio de ese apellido, bajo esa intendencia. 2014, ampliación del objeto social a la comercialización y adquisición de inmuebles. 2018, designación de Héctor Miguel Calabró\index{Calabró, Héctor Miguel} como gerente suplente. 2021, aprobación del plano del loteo. Cada uno de esos actos es público y ninguno, por sí mismo, es objetable.

El vínculo familiar, si existe, se acredita con partidas del Registro Civil, y es uno de los puntos que el pedido de derecho a réplica debe plantear de manera directa antes de que este pasaje se publique.
\end{cautela}

\subsection*{Dos hallazgos laterales de las actas}

Las actas de proclamación son documentos exhaustivos --- consignan cada cargo de cada municipio de la provincia con nombre y documento --- y de ellas se desprenden dos datos que este libro no buscaba.

\begin{cautela}
\textbf{El primero es un apellido a caballo del límite.} El acta de 2015 registra, para el municipio de Vaqueros, al intendente \textbf{Daniel Roberto Moreno Ovalle\index{Moreno Ovalle, Daniel Roberto}}. Y registra, para el municipio de \textbf{Salta capital}, al concejal \textbf{Mario Enrique Moreno Ovalle}, electo por la misma fuerza --- Partido de la Victoria --- en la misma elección.

Los números de documento de ambos, que el acta consigna y este libro no reproduce, son \textbf{consecutivos}. No se afirma parentesco, que se acredita con partidas y no con numeración, pero deja constancia de que dos personas de igual apellido compuesto y documento de numeración correlativa ocuparon simultáneamente, en 2015, un cargo ejecutivo en Vaqueros y uno deliberativo en la capital.

Y eso importa por lo que el capítulo \ref{cap:vaqueros} estableció: el río Vaqueros es el \textbf{límite} entre ambas jurisdicciones, y ninguna de las dos tiene competencia sobre el conjunto del cauce. Que un mismo apellido ocupara posiciones en las dos orillas en el mismo período es un dato de estructura política, no una imputación, y merece verificarse antes que interpretarse.
\end{cautela}

\textbf{El segundo es un apellido antes de tiempo.} Entre los concejales electos de La Caldera en 2015 figura \textbf{Alfio Edgardo Sumbay}, por el Frente Renovador. Cuatro años después, Diego Ramiro Sumbay\index{Sumbay, Diego Ramiro} ganaría la intendencia por una fuerza distinta, el Movimiento Comunitario Pluricultural. El apellido está en el departamento mucho antes: en 1965 un Jenaro Sumbay promueve la posesión treintañal de 115 hectáreas en el partido de Yacones (edicto 20712, B.O.\ Nº 7366, h.\ 15); parentesco, ninguno que conste.

De nuevo: mismo apellido, sin parentesco acreditado. Lo que la observación aporta al capítulo es que el apellido que hoy conduce el municipio ya tenía representación en el Concejo cuatro años antes, y por otra fuerza política. La política departamental no se lee sólo por partidos.

Un tercer dato, menor pero útil para el capítulo \ref{cap:poblacion}. En 2011 y en 2015, tanto La Caldera como Vaqueros eligieron \textbf{tres concejales} por acta. Si el Concejo tenía cinco miembros, como los de segunda categoría de la Ley 1349 entonces vigente, eso indicaría renovación parcial del cuerpo; las actas no dicen cuántos lo integraban. Y Moreno Ovalle\index{Moreno Ovalle, Daniel Roberto} fue electo intendente de Vaqueros en 2011 por Salta Somos Todos y en 2015 por el Partido de la Victoria: dos sellos distintos en dos elecciones consecutivas, antes de los que usaría en 2019 y 2023.

\textbf{De dónde venía el error.} El listado no corroborado proponía para 1983 a \textbf{Roberto Adán Galli}. Ese nombre existe y corresponde a un intendente real, pero de otra jurisdicción: Galli fue el \textbf{primer intendente de la Ciudad de Salta} tras el retorno democrático, sucediendo al interventor Néstor Salvador Quintana el 10 de diciembre de 1983 (artículo «Néstor Salvador Quintana» de Wikipedia en español, fuente terciaria, consultado el 18 de septiembre de 2026; y la necrológica de Galli en \textit{Iruya.com}, 26 de mayo de 2013, \url{http://noticias.iruya.com/newnex/sociedad/gente/7774-ha-fallecido-en-salta-el-ingeniero-roberto-adan-galli.html}). Y no es el único: \textbf{Alberto Javier Alderete}, que el mismo listado proponía para 1987, fue según \textit{La Gaceta Salta} (16 de mayo de 2015) el primer intendente de la capital elegido por voto directo. El listado había tomado la serie de la capital y la había atribuido al departamento, del mismo modo que había fundido a Mendaña\index{Mendaña, Luis Gerardo} con Escalera\index{Escalera, Daniel Alejandro}.

Se consigna porque ilustra el riesgo con precisión: un resumen de segunda mano no inventa de la nada, mezcla series contiguas. Y el resultado tiene la apariencia exacta de un dato.

\begin{cautela}
\textbf{Y hay una razón por la que ese tramo no está en las actas del Tribunal Electoral.}

Galli, según esas dos fuentes ---ninguna de ellas un acto administrativo: el decreto de designación sigue sin obtenerse---, no fue electo sino \textbf{designado por el gobernador} Roberto Romero. Eso es coherente con el artículo 22 de la Ley 1349 --- el que el apéndice \ref{cap:normativa} señala como fósil administrativo ---, según el cual los intendentes ``serán nombrados y removidos por el P.E.\ con acuerdo del Senado''. Ese régimen rigió hasta que la Constitución provincial de 1986 introdujo la elección directa.

De donde se sigue que, entre 1983 y 1987, \textbf{es probable que no existan actas de proclamación de intendentes electos}, porque los intendentes no se elegían: se designaban. Buscarlas en el Tribunal Electoral sería buscar en el lugar equivocado. Lo que habría para ese período son \textbf{decretos de designación} del Poder Ejecutivo provincial, con acuerdo del Senado, publicados en el Boletín Oficial.

Para La Caldera y Vaqueros, que eran Comisiones Municipales (véase más arriba), lo que habría son decretos de designación de \textit{presidentes de comisión}, que el artículo 33 dejaba en manos del Poder Ejecutivo sin acuerdo del Senado; el acuerdo regía para los intendentes de primera y segunda categoría, como el de la capital. Para 1983, la Ley 6131 y la Ordenanza 047 ya lo confirman en La Caldera, y el decreto que designó a su presidente es el 877 del 1 de septiembre de 1982 (véase más arriba); de ahí la consecuencia: para 1984--1987 hay que buscar decretos, no proclamaciones. Para 1984 y 1985 la lectura los encontró: el 145 del 16 de enero de 1984 y el 2862 del 28 de diciembre, que designa por un nuevo período al mismo presidente (véase más arriba). Y para 1986 y 1987 también: el 229 de febrero de 1986 lo vuelve a designar, y el 13 de mayo de 1987 el 907 le acepta la renuncia y el 908 designa a su sucesor; en junio la convocatoria para el 6 de septiembre pone al intendente en la boleta (véase más arriba).
\end{cautela}


